% concatenated from sprouts.tex
\documentclass[11pt,reqno]{amsart}
\usepackage{amsmath, amsthm, amsfonts, amssymb}
\usepackage{graphicx, amscd, color,bm}    
  %\usepackage[figuresonly,nolists,nomarkers]{endfloat}
%\renewcommand{\processdelayedfloats}{}
\usepackage[colorlinks]{hyperref}
\usepackage[T2A]{fontenc}
\usepackage[utf8]{inputenc}
\usepackage[mathscr]{eucal}


\newtheorem{theorem}{Theorem}[section]

\newtheorem{corollary}[theorem]{Corollary}

\newtheorem{definition}[theorem]{Definition}
\newtheorem{example}[theorem]{Example}
\newtheorem{lemma}[theorem]{Lemma}
\newtheorem{notation}[theorem]{Notation}
\newtheorem{proposition}[theorem]{Proposition}
\newtheorem{remark}[theorem]{Remark}
\newtheorem{summary}[theorem]{Summary}
\def \odeg {\mathop{\rm odeg}\nolimits}
\newcommand{\be}{{\bm e}}
\newcommand{\Ord}{\mathop{\rm Ord}\nolimits}
\newcommand{\bmb}{{\bm \beta}}
\newcommand{\bi}{{\bm i}}
\newcommand{\bj}{{\bm j}}
\newcommand{\IN}{{\subseteq}}
\newcommand{\NI}{{\supseteq}}
\newcommand{\ga}{{\gamma}}
\newcommand{\Ga}{{\Gamma}}
\newcommand{\al}{{\alpha}}
\newcommand{\hga}{{\hat\gamma}}
\newcommand{\Nfi}{{N_\phi}}
\newcommand{\eS}{{\mathcal S}}
\begin{document}
\title{Self-similar dendrites with finite boundary and $P$-sprouts}
\thanks{The study was carried out under the state contract of the Sobolev Institute of Mathematics (project FWNF-2026-0026).}
\date{}
\author[Andrei Tetenov, Ivan Yudin, Dmitrii Drozdov]{Andrei Tetenov, Ivan Yudin, Dmitrii Drozdov}
\begin{abstract}
Each self-similar dendrite $K$ with a finite self-similar boundary defines a finite acyclic edge-labeled bipartite graph $\Gamma$,  called the sprout of $K$. The paper shows that the sprout  $\Gamma$   determines the combinatorial properties of the dendrite $K$ and its topological structure.\\


\end{abstract}
\maketitle
\textbf{Mathematics Subject Classification.} Primary: 28A80, 54F50; Secondary: 05C20.\\
\textbf{Keywords and phrases:} self-similar set,  dendrite,   intersection graph,  single intersection property, intersection graph, ramification points, acyclic fractal set.




\section{Introduction}
Let $\eS=\{S_1,...,S_m\} $ be a system  of contractive injections of a complete metric space $X$.
A  compact non-empty set $K$ is called the attractor of a system $\mathcal{S}$ if it satisfies the equation
$K=S_1(K)\cup...\cup S_m(K)$.   The existence and uniqueness of the attractor $K$ was proved by J.Hutchinson \cite{Hut}. 
The subsets $K_{i_1...i_n}= S_{i_1}...S_{i_n}(K)$ are called  copies of $K$. 
The set of those points $x\in K$, whose images are the intersection points of different copies of $K$, is the self-similar boundary of $K$. 
If the attractor $K$ is connected and does not contain a closed curve, it  is called a self-similar dendrite. 
The main object of this paper is self-similar dendrites with a finite boundary. 
These dendrites possess the single intersection property (SIP) and therefore have a clear combinatorial characterization.\\

The study of dendrites dates back to papers \cite{Maz, Men, Mil, W} that appeared in the early 1920-s and defined key objectives in dendrite topology. The main work of that period is the doctoral dissertation \cite{W} of Tadeusz Vazhevsky. He gave a rigorous systematic description of the topology of dendrites and proved the existence of a unique locally connected continuum with the property of universality. In such a space, the ramification points have a strictly defined order, and any other dendrite with a similar or lesser ramification order can be homeomorphically embedded in it. Karl Menger was one of the main researchers of the theory of curves and dendrites in the 1920-s. In \cite{Men} he developed a general theory of curves. He also gave a classification of dendrite points and investigated their properties in the context of dimension theory. S.Mazurkevich laid the foundations for the study of locally connected continua in \cite{Maz}. In subsequent articles  he investigated the conditions under which a continuum is a dendrite and dealt with the issues of their embeddability. K. Zarankiewicz formulated the topological question of whether dendrites are "non-contractible", that is, not homeomorphic to their own subset. In 1932 \cite{Mil}  Miller refuted this hypothesis by constructing a counterexample known as the Miller dendrite, which is homeomorphic for its part.
An exhaustive overview of the research in this area  can be found in \cite{Char}. 
The inception of the theory of self-similar sets \cite{Hut} initiated the study of self-similar dendrites. In  1985, M. Hata  \cite{Hata85} proved that the set of endpoints of a non-trivial self-similar dendrite is infinite. In 1995, J. Kigami \cite{Kig,Kig95} investigated the shortest-path metric in post-critically finite self-similar dendrites and constructed regular Dirichlet forms for such dendrites. Under certain conditions, Julia sets of complex polynomials may be dendrites. 
Hubbard and Douady \cite{DH1984} described this case in detail. Such Julia sets arise if the critical point is strictly preperiodic. Hubbard trees were introduced to study their structure. This is a finite tree that connects all the points of the critical point orbit inside the Julia set. It contains all the combinatorial information about the Julia set, allowing the classification of polynomials.
As  follows from the results of A.Kameyama \cite{Kam}, if a self-similar dendrite is the attractor of some invertible system of contractions, it is homeomorphic to a Julia set of some polynomial.
 The article \cite{SSS2} by C. Bandt and K. Keller deserves special attention, as it contains an intersection graph criterion for a  self-similar continuum with single intersection property to be a dendrite (Proposition 9). The preprint \cite{SSS6} by C. Bandt and J. Stanke introduces the idea of a 'main tree' for post-critically finite self-similar dendrites.
The topology automaton approach to the study of self-similar dendrites was considered in \cite{Bandt2024, RZ2020, RWYZ2023}.
Self-similar dendrites with non-trivial intersection patterns and infinite ramification were studied in \cite{AST2024, AK2025}.
The growth of self-similar dendrites was considered in \cite{barua2022}.

A simple and intuitive way to construct self-similar dendrites in the plane and in space is provided by polygonal (respectively, polyhedral) systems.
This method was described in \cite{STV} and expanded in \cite{Adam}. 


In this paper, we develop a method for defining self-similar dendrites with single point intersections using simple combinatorial schemes called sprouts. These schemes determine self-similar dendrites up to isomorphism. We provide an algorithm that enables the reconstruction of the main topological elements of a self-similar dendrite from its sprout.


The main theorems proved in the paper are as follows.

The isomorphism theorem for self-similar dendrites (Theorem \ref{isom}) states that  if systems ${\eS}=\{S_1,...,S_m\}$ and $\tilde {\eS}=\{\tilde S_1,...,\tilde S_m\}$ have isomorphic $P$-sprouts $\Gamma$ and $\tilde \Gamma$, then their attractors $K$ and $ \tilde K$ are isomorphic. 

Theorem \ref{inf_ram} states the conditions under which  the  set of addresses of a boundary point $p\in \partial K$ is finite, countable, or uncountable and gives a formula for finding  ${\rm Ord}(p,K)$.


The third is  Theorem \ref{ord_b}   on the order of boundary points with respect to the main tree. \\ 

In Section 3 we introduce a $P$-sprout $\Gamma$ that defines a self-similar dendrite $K$ with a finite self-similar boundary and  define the index diagram of $\Gamma$.

In Section 4 we consider   SIP systems of continua and refinement sequences of SIP systems, which serve as a tool to prove  Theorem \ref{isom}).

In Section 5 we consider the main tree of a self-similar dendrite, evaluate the orders of the points of the attractor, and determine the cardinality of the set of addresses of  boundary points. The main statements of this section are Proposition \ref{numwalks} and Theorem \ref{inf_ram}.

In section 6 we consider a semigroup $G_\phi $ of  transformations of the set $P$, and prove  Theorem \ref{ord_b} on the orders of boundary points and ramification points of the main tree. Theorem \ref{radr}
gives a method of finding the addresses of the ramification points of the main tree.

\section{Preliminaries}

{\bf Self-similar sets.} Let ${\eS}=\{S_1,  S_2, \ldots,  S_m\}$ be a system of injective contractions in a complete metric space $(X,  d)$. 
There is a unique compact non-empty set  $K{\subset} X$ such that $K =\bigcup\limits_{i = 1}^m S_i (K)$. The set $K=K({\eS})$  is called the attractor of the system  ${\eS}$ \cite{Hut}. The subsets $K_i=S_i(K), i\in\{1,...,m\}$, are called {\em copies} of the attractor $K$. 


The set $I=\{1, 2, \ldots, m\}$ is called the {\em set of indices},  the set $I^*=\bigcup\limits_{n=1}^\infty I^n$ 
is the set of  {\em multiindices} ${\bf {j}}=j_1j_2\ldots j_n$. For a multiindex
${\bf {j}}$, the map $S_{\bf {j}}=S_{j_1}\cdot S_{j_2}\cdot\ldots\cdot S_{j_n}$ defines a copy $K_{\bf {j}}=S_{\bf {j}}(K)$ of order $n$. The set $G_{\eS}=\{S_{\bf {i}},{\bf {i}}\in I^*\}$ is a semigroup, generated by the system ${\eS}$.\\

Since any sequence 
${\bf\alpha}=\alpha_1\alpha_2\ldots\alpha_n\ldots\in I^\infty$ specifies a unique point 
$\pi({\bf\alpha})=\bigcap\limits_{n=1}^\infty K_{\alpha_1\ldots\alpha_n}$,
we define a map
$\pi:I^{\infty}\rightarrow K$  called the {\em index map}.  If $\pi({\bf\alpha})=x$,  then ${\bf\alpha}$ is called {\em an address} of the point $x$.  
 We say that an address ${\bf\alpha}$ is preperiodic if ${\bf\alpha}={\bf {j}}\bar{\bf {i}}$ for some ${\bf {i}},{\bf {j}}\in I$.

Given an address ${\bf\alpha}=\alpha_1\alpha_2\ldots\alpha_n\ldots$ of a point $y\in K$ we consider a sequence of points $y_k=\pi(\alpha_{k+1}\alpha_{k+2}\ldots)= S_{\alpha_1\alpha_2\ldots\alpha_k}^{-1}(y)$ which are called {\em the predecessors} of the point $y$.
% A point $y_1 \in K$ is a {\em (first) predecessor} of a point $y \in K$ if there is  $i \in I$ such that $y = S_{i}(y_1)$.
In other words, a point $y_k \in K$ is a {\em $k$-th predecessor} of the point $y \in K$ if there is a word ${\bf {i}} \in I^k$ such that $y = S_{{\bf {i}}}(y_k)$.  \\



The union ${\mathcal C}=\bigcup\limits_{i\neq j} K_i\cap K_j$ is called { \em the critical set} of ${\eS}$. 
The set  of all predecessors of points of ${\mathcal C}$, i.e. the set of all  $x\in K$ such that for some ${\bf {i}}=i_1\ldots i_n$, $S_{{\bf {i}}}(x)\in {\mathcal C}$, is called {\em the self-similar boundary} ${\partial} K$ of the attractor $K({\eS})$ \cite{Mor}.
For any $n\in \mathbb{N}$ we define ${\mathcal C}^{(n)} = \bigcup\limits_{|\bf{i}| = n-1} S_{\bf{i}} ({\mathcal C})$.

Two multiindices  ${\bf {i}},{\bf {j}}\in I^*$ are incomparable, if neither of them is an initial subword of the other.

For any incomparable multiindices ${\bf {i}},{\bf {j}}\in I^*$, $K_{\bf {i}}\cap K_{\bf {j}}=S_{\bf {i}}({\partial} K)\cap S_{\bf {j}}({\partial} K)$. \\
The self-similar set $K({\eS})$ {\em has a finite self-similar boundary} if the set ${\partial} K$ is finite. 
The set $K({\eS})$ has the {\em finite intersection property} (FIP) if all the sets $K_i\cap K_j$ are finite and it has the {\em single intersection property}   (SIP) if for any $i,j\in I$, $\#(K_i\cap K_j)\le 1$.\\


{\bf Dendrites and sprouts.} We call the attractor $K$  a {\em self-similar continuum} if it is connected. 
By Hata's Theorem \cite{Hata85}, any self-similar continuum is locally connected.\\
A {\em  dendrite}  is a locally connected continuum $K$ that does not contain a simple closed curve \cite{Kur}. 
Therefore, each acyclic self-similar continuum is a self-similar dendrite.\\

  \begin{definition}\label{minarc}
      
Let $D$ be a dendrite, $A, B$ be connected disjoint subsets of $D$. Then there is a single arc $\gamma(x,y)$ connecting the points $x\in A$ and $y\in B$ such that $\gamma \cap A =\{x\}, \gamma\cap B=\{y\}$. We will call such an arc $\gamma$ the minimal arc connecting $A$ and $B$.
\end{definition}
 




%If the set $\pi^{-1}(\partial K)$ of addresses of  boundary points is finite, then the 
% system ${\eS}$ is called postcritically finite.
If $K$ is a   self-similar dendrite, then  each non-empty intersection   $P_{ij} =S_i(K)\cap S_j(K)$, where $ i \neq j$, is a singleton $\{p\}$. 
Thus,  $K$ has the single intersection property \cite{TYK}. 
Consequently, the set $\bigcup\limits_{i\in I} S_i(\partial K)$ subdivides $K$ into  pieces whose closures are  copies $K_i$.\\




We consider the following two sets: 
 The set $W = \{K_i\}_{i \in I}$ of copies of $K$ and   the set $B = {\mathcal C} \cup \partial K$ of intersection points of copies of $K$ and of boundary points of $K$.\\


\begin{figure}[htp] 
\centering
\includegraphics[width=\textwidth]{pictures/trapezium2.pdf}
{\Large$4 \mathop{\longrightarrow }\limits^{w_1}\ 1  \mathop{\longrightarrow }\limits^{w_1}\ 2  \mathop{\longrightarrow }\limits^{w_2}\ 1 \mathop{\longrightarrow }\limits^{w_1} 2\mathop{\longrightarrow }...$}
\caption{The attractor $K$ is on the left, its $P$-sprout is at center  and index diagram is on the right. The sequence of predecessors of $p_4$ is shown below.}
\label{f1}
\end{figure}

If the system ${\eS}$ has the finite intersection property, we define its {\em intersection graph} $\Gamma({\eS})$ \cite{TYK} as
an edge-labeled bipartite graph $(W, B; E)$ with parts $W$ and $B$,  in which an edge $ e = (K_i, p ) \in E$ 
if and only if $p \in K_i$. Here, a label assigned to the edge $e=(K_i,p)$  corresponds to the boundary point $p_j \in {\partial} K$, for which $S_i(p_j)=p$. 
This labeled intersection graph $\Gamma({\eS})$ will be called {\em the sprout} of the system ${\eS}$.\\ 



Together with the graph $\Gamma$, we consider an edge-labeled directed graph ${{\mathcal G}}$, derived from $\Gamma$, which we  call {\em the index diagram} of the system ${\eS}$.  
It shows  the relationship between the boundary points of the attractor $K({\eS})$.
The set of vertices of this graph is ${\partial} K$. 
For $x_i,x_j\in {\partial} K$, there is a directed edge from  $x_i$ to $x_j$ with edge label $w_k$ iff there is  $k\in I$ such that $S_k(x_j)=x_i$. \\

The attractors $K({\eS})$, $K({\EuScript T})$  of  systems of injective contractions ${\eS}=\{S_1,\ldots,S_m\}$ and ${\EuScript T}=\{T_1,\ldots,T_m\}$   are isomorphic if there is a homeomorphism $\varphi:K({\eS}) \to K({\EuScript T})$ such that for any $x\in K({\eS})$ and $i\in I$, $\varphi(S_i(x))=T_i(\varphi(x))$. 
If the latter is fulfilled, the sprouts $\Gamma({\eS})$ and $\Gamma({\EuScript T})$ are also    equivalent.
On the other hand, we show  that if  the sprouts $\Gamma({\eS})$ and $\Gamma({\EuScript T})$ are equivalent, then the  attractors $K({\eS})$ and $K({\EuScript T})$ are isomorphic.\\



Fig.\ref{f1}  shows an example of a self-similar set $K({\eS})$, where ${\eS}=\{S_1, S_2\}$, the sprout $\Gamma({\eS})$ and the index diagram ${{\mathcal G}}({\eS})$.  One sees that the self-similar boundary ${\partial} K$ consists of 4 points, $p_1,p_2,p_3,p_4$ shown in red numbers. 

The only critical point in $K$ is  $S_1(p_4)=S_2(p_3)=p_5$. 
All of this  gives us the sprout shown in the center.
The arrows at the edges of the index diagram on the right are directed toward the predecessors of the boundary points. Thus, each directed walk in the index diagram gives a sequence of predecessors of the initial point of the walk. The sequence of labels in this walk yields the address of the initial point. The sequence below the three pictures  shows the  infinite walk starting from $p_4$. At the same time, the numbers on the labels form the address 
$112121212...$ of  the point $p_4$.





The existence of a directed walk $w_{i_1}...w_{i_n}$ from $a$ to $b$ in this diagram implies that
$S_{i_1...i_n}
(b) = a$ and  $a \in K_{i_1...i_n}$. Thus, each infinite directed walk starting from the point $a$ defines an address of the point $a$. The number of addresses of a boundary point $a$ is equal to the number of infinite walks in
the index diagram that start from $a$, as shown on the right.






  
\section{Definition of a  $P$-sprout}

\begin{definition}
Let $\Gamma=(B,W,E)$ be an acyclic bipartite  graph with parts $(B,W)$.
Let $P\subset B$; let $\varphi:E\to P$ be a  surjection such that for any $w\in W$, the restriction of $\varphi$ to the subset $E(w)\subset E$ is injective. 
Then the pair $(\Gamma,\varphi)$, is called a $P$-sprout.
\end{definition}


The {\em critical set} of a sprout $(\Gamma,\varphi)$ is the set $C=\{b\in B:{\mathop{\rm Ord} \nolimits}(b,\Gamma)>1\}$. 

For  $b\in B$, we denote by ${\widetilde\varphi}(b)$ the set $\varphi(E(b)$) of all $p\in P$ corresponding to the edges incident with $b$. {\em The boundary of the sprout} $(\Gamma,\varphi)$ is the set
${\partial}\Gamma=\bigcup\limits_{n=1}^\infty \widetilde\varphi^n(C)$.



The sprout $(\Gamma,\varphi)$ is {\em correctly defined} if its boundary ${\partial}\Gamma$ is equal to $P$. (see Fig.2 for incorrectly defined sprout). In the following, it will be assumed that all sprouts are correctly defined.

 Moreover, if for any vertex $b\in B\setminus P$ and for any vertex $w\in W$, its degree in $\Gamma$ is greater than 1, then    
$(\Gamma,\varphi)$ is called a {\em regular} $P$-sprout.\\

{\bf Remark:} In the regular $P$-sprout each component of the complement to the vertex of the sprout contains at least one boundary point.


\begin{figure}[htp]
\centering    
\includegraphics[width=\textwidth]{pictures/m1_5.pdf}
\caption{ A sprout on the left is not  correctly defined ($3 
    \notin \bigcup\limits_{n=1}^{\infty} \tilde{\varphi}^n(C)
    $), a sprout on the right is defined correctly  but is not regular.  }
\label{fig:cor_def}
\end{figure}


Two sprouts $(\Gamma(B,W,E),\varphi:E\to P)$, $(\Gamma'(B',W',E'),\varphi':E'\to P')$ are {\em  isomorphic} if there is a bijection $x\to x'$ between the respective  sets mentioned above that preserves the incidence relation and such that for any $e\in E$, $(\varphi(e))'=\varphi'(e')$.\\





Let $E_P$ be a  set of edges of $\Gamma$ incident with the vertices $p\in P$.


 

\begin{definition}
The index diagram of the sprout $(\Gamma,\varphi)$ is  a digraph ${{\mathcal G}}_P=(P,{{\mathcal E}}, \hat\varphi)$  whose vertex set is $P$.  
Each edge $e=(w_k,p_i)\in E_P$ with its label $\varphi(e)=p_j\in P$ defines a directed edge $\vec{e} = ({p_i,p_j}) \in {{\mathcal E}}$ with the label $\hat\varphi(\vec e)=w_k$.
\end{definition}


The paper \cite{Adam} considers the deformations of a polygonal system ${\eS}$ that preserve the sprout of the system ${\eS}$ and produce a parametrized family of isomorphic self-similar dendrites. 

% For any $p_i \in P$ we define the set
  %  $\tau(p_i) = \{ p_j:\; ({p_i,p_j}) \in {{\mathcal E}} \}$
 %   and $\omega(p_i) =  \{ p_j: ({p_j,p_i}) \in {{\mathcal E}} \}$. 
 %   Consequently, the set of all outgoing edges from the vertex $p_i$ will be denoted by ${{\mathcal E}}_\tau(p_i) = \{ (p_i,p_j): p_j\in \tau(p_i) \}$, and the set of all incoming edges to the vertex $p_i$ will be denoted by ${{\mathcal E}}_\omega(p_i) = \{ (p_j,p_i): p_j\in \omega(p_i) \}$.

Since the graph $\Gamma$ is a tree, for any $p\in P$ and $w\in W$ there is at most one edge $e$ in $\Gamma$ incident with $p$ and $w$. Let $\varphi(e)=p'$. 
Then there is at most one edge $\vec e=(p,p')\in{{\mathcal E}}$ in ${{\mathcal G}}_P$ such that $\hat\varphi(\vec e)=w.$
Therefore, the  restriction of $\hat\varphi$ to the set of outgoing edges from  $p_i$ is injective.

On the other hand, the existence of two incoming edges $\vec {e_1}=(p_1,p)$,  $\vec {e_2}=(p_2,p)$ such that $\hat\varphi(\vec {e_1})=\hat\varphi(\vec {e_2})=w$ contradicts the injectivity of $\varphi|_{E(w)}$ in $\Gamma$. 
Therefore, the  restriction of $\hat\varphi$ to the set of edges of ${{\mathcal E}}$,  incoming  to  $p_i$, is injective.\\

{\bf The sprout, defined by an IFS.} If a system ${\eS}=\{S_1,...,S_m\}$ with the attractor $K$ and a finite boundary $\partial K=P$ has the single intersection property, its $P$-sprout $\Gamma({\eS})$ is the edge-labeled bipartite intersection graph $\Gamma(W, B, E,\varphi)$ with parts $W=\{K_1,...,K_m \}$ and $B={\mathcal C}\cup \partial K$, in which an edge $ e = (K_i, p ) \in E$ 
iff $p \in K_i$. 
A label $\varphi(e)$ assigned to  $e=(K_i,p)$, is the unique point $p_j \in {\partial} K$ for which $S_i(p_j)=p$. 

\begin{definition}\label{walks}
Let ${{\mathcal G}} =(P,{{\mathcal E}})$ be a   directed graph. A sequence of vertices and edges $\omega(p_0,p_k)=(p_0,\vec{e_1},p_1,\hdots \vec{e_k},p_k)$ is {\em a walk from $p_0$ to $ p_k$ } in ${{\mathcal G}}$ if  for any $i=1,...,k$,  $\vec e_i=(p_{i-1},p_i)$. 
If $p_0=p_k$, then $\omega$ is a {\em cycle} and its vertices are {\em cyclic vertices}.\\
A walk $\omega$ starting from a point $p$ is denoted by $\omega(p,...)$ and a walk that has an endpoint $q$ is denoted by $\omega(...,q)$.
We write $\omega(p,\sigma)$ for a walk whose infinite subwalk is contained in $\sigma$.
\end{definition}

The set of all   walks $\omega(p,q)$ in ${{\mathcal G}}$  is denoted by $\Omega(p,q)$. 
The set of all infinite walks $\omega(p,..)$ in  ${{\mathcal G}}$ starting from a vertex $p\in P$ is denoted by $\Omega(p)$. 
The set of all infinite walks $\omega(p,\sigma)$ in  ${{\mathcal G}} $ is denoted by $\Omega(p,\sigma)$.

If needed, we may represent a walk $\omega$ as a sequence  $(p_0,p_1,\hdots, p_k)$ of its vertices, or a 
sequence $( \vec{e_1},\hdots \vec{e_k})$  of its edges. \\

Let  the attractor $K({\eS})$ of a system ${\eS}=\{S_1,  S_2, \ldots,  S_m\}$  be a self-similar dendrite with a finite   boundary $P$,  $\Gamma({\eS})$ be its sprout, and ${{\mathcal G}}_P$ its index diagram. 
Let $\omega(p_0,p_k)=(p_0,\vec{e_1},\hdots \vec{e_k},p_k)$ be a walk  in ${{\mathcal G}}_P$  and $\hat\varphi(\vec {e_j})=w_{i_j}$. 
Then $p_0=S_{i_1i_2...i_k}(p_k)$. 
We call $\hat\varphi(\omega)=i_1i_2...i_k$ a multiindex defined by the walk $\omega$. Similarly,  an infinite walk $\omega(p_0,...)=(p_0,\vec{e_1},p_1, \vec{e_2}....)$ in ${{\mathcal G}}_P$ defines an infinite string $\hat\varphi(\omega)=\alpha=i_1i_2....$.

\begin{proposition}\label{infwalks}
Let $\omega(p_0,...)$ be an infinite walk in ${{\mathcal G}}_P$. The infinite string $\hat\varphi(\omega)=i_1i_2....$ defined by the walk $\omega$ is the address of the point $p_0$ in $K$ and $\pi^{-1}(p_0)=\hat\varphi( \Omega(p_0))$.
\end{proposition}

\begin{proof}
Any initial subwalk $\omega_k=\omega(p_0, p_k)$ of $\omega(p_0,...)$ defines a copy $S_{i_1...i_k}(K)$ that contains $p_0$; these copies form a nested sequence $K_{i_1}{\supset}...{\supset} K_{i_1...i_k}{\supset}...$, whose intersection is $\{p_0\}$. 
Consequently, $\hat\varphi(\omega(p_0,...))=i_1i_2...$ is the address of the vertex $p_0$.    
\end{proof}

Each address $\alpha=i_1i_2...$ defines a unique point $x=\pi(\alpha)$ in a self-similar set $K$. This imposes the following restriction on the sprouts under consideration.

\begin{definition}
We say that $P$-sprout $\Gamma$ is {\em admissible} if and only if there are no two different boundary points that have the same address. 
\end{definition}
\begin{figure}[htp]
\begin{center}
\includegraphics[width=.7\textwidth]{pictures/unf.pdf}
\caption{Example of an inadmissible $P$-sprout (on the left) and its index diagram (on the right).}
\label{fig:unf}
\end{center}
\end{figure}
\begin{lemma}\label{ordpi}
Let  the attractor $K({\eS})$ of a system ${\eS}=\{S_1,  S_2, \ldots,  S_m\}$  be a self-similar dendrite with  a finite   boundary.
For any $x\in K$, ${\mathop{\rm Ord} \nolimits}(x,K)\ge\#\pi^{-1}(x)$ if the set $\pi^{-1}(x)$ is finite and ${\mathop{\rm Ord} \nolimits}(x,K)$ is infinite if $\pi^{-1}(x)$ is infinite.
\end{lemma}

\begin{proof}
If the point $x\in K$ has $l$  different  addresses  $\alpha^k=i^k_1i^k_2....$, then there is $n\in\mathbb{N}$ such that all multiindices ${\bf {i}}^k=i^k_1i^k_2....i^k_n$   are different. 
Consider  two copies $K_{\bf {i}^j}$ and $K_{\bf {i}^k}$. 
Their intersection is $\{x\}$. Consequently, for any connected components $U,V$ of the set $K\setminus\{x\}$ such that $U\cap K_{\bf {i}^j}\neq\varnothing$ and $V\cap K_{\bf {i}^k}\neq\varnothing$, $U\cap V=\varnothing$. Therefore, the cardinality of the set of components of the set $K\setminus\{x\}$ is greater than or equal to $\#\pi^{-1}(x)$.
\end{proof}

\section{Isomorphism theorem  for the attractors.}

 All  definitions and statements of this section will not be used in the next sections,
 so the reader may feel free to omit it.

\begin{definition} 
Let $\Gamma(W,B,E,\varphi)$ be a $P$-sprout. Put $W'=W\times W$.\\ Put $B'=(W\times B\cup B)/\sim$, where $w\times  b_1\sim b_2$ if $b_2\in P, (w,b_1)\in E$ and $\varphi(w,b_1)=b_2$. Put $E'=W\times E$ and $\varphi'(w\times e)=\varphi(e)$.\\
The $P$-sprout $\Gamma(W',B',E', \varphi')$ is denoted by
 $\Gamma^2$.(see Fig.\ref{G^2})\\  
The index diagram ${{\mathcal G}}^2_P$ for the sprout $\Gamma^2$ has the same set $P$ of vertices; the edges are all  walks $\vec{e_i} \vec{e_j}$ of length $2$ in ${{\mathcal G}}_P$.
\end{definition}



\begin{lemma}
Let  the attractor $K({\eS})$ of a system ${\eS}=\{S_i,i\in I\}$  be a self-similar dendrite with  a finite   boundary ${\partial} K=P$ and let $\Gamma$ be its $P$-sprout. \\ The system $\{K_{ij},i,j\in I\}$ has the SIP and is the refinement of the system $\{K_i, i\in I\}$. (see Def. \ref{refin})\\ 
 The sprout $\Gamma^2$ is the $P$-sprout of the system ${\eS}^2=\{S_iS_j,(i,j)\in I^2\}$ and $\Gamma^2$ is the labeled intersection graph for the system of copies $K_{ij}$ of $K$. 
\qed
\end{lemma}

\begin{figure}[htp]\label{G^2}
\begin{center}
\includegraphics[width=\textwidth]{pictures/g2a.pdf}
\caption{\small The left column  shows how the sprout $\Gamma^2$(at the bottom) is obtained from the sprout $\Gamma$ (at the top). The  middle picture shows the neighborhoods of points 1,2 and 3 that shrink to each of these points. Right column shows the systems of  copies $K_i$ and $K_{ij}$. }
\end{center}
\end{figure}

\begin{corollary}
If systems ${\eS}=\{S_1,...,S_m\}$ and $\tilde {\eS}=\{\tilde S_1,...,\tilde S_m\}$ have isomorphic sprouts $\Gamma$ and $\tilde \Gamma$, then for any $n\in \mathbb{N}$, the sprouts $\Gamma^{2^n}$ and $\tilde{\Gamma}^{2^n}$ are isomorphic.
\qed
\end{corollary}


% \begin{figure}[htp]
% \begin{center}
% \includegraphics[width=.7\textwidth]{exABC.pdf}
% \includegraphics[width=.7\textwidth]{index.pdf}
% \caption{The index diagrams of a 3-sprout, a 4-sprout, and a 3-sprout. The elements $p_i$ are marked by the red numbers, the values of $\varphi(e)$ are marked in black   above the edges.
% }
% \label{fig:exABC}
% \end{center}
% \end{figure}

% \begin{definition}
% The digraph  ${{\mathcal G}}_I=(I,{{\mathcal E}})$ is called the index diagram of an m-sprout  $\Gamma$
% \end{definition}


% \begin{figure}[htp]
% \begin{center}
% \includegraphics[width=.8\textwidth]{pictures/sproutA.pdf}
% \caption{$P$-sprout $\Gamma$ on the left and its index diagram on the right.
% }
% \label{fig:sproutA}
% \end{center}
% \end{figure}




\begin{definition}\label{fipss}\cite{TYK}
Let $\mathcal K=\{K_i, i\in I=\{1, \ldots, m\}\}$ be a finite system of continua in a complete metric space $(X,d)$.  We say $\mathcal K$ has the {\em single intersection  property (SIP)},  if for any  $i\neq j\in I $, the set $P_{ij}=(K_i\cap K_j)$ is empty or a singleton.  
\end{definition}

In the setting of Definition \ref{fipss} we denote $|\mathcal K|=\bigcup\limits_{i\in I}K_i$,   $P=\bigcup\limits_{i\neq j}P_{ij} $, and $P_i=\bigcup\limits_{j\in I\setminus\{i\}}P_{ij}$.  Regarding $|\mathcal K|$ as the subspace of $X$ supplied by the induced topology, we see that for any $i$ the set $P_i$ is the boundary $\partial K_i$ of the set $K_i$ in $|\mathcal K|$,  and that its interior is $\dot K_i=K_i\setminus P_i$.  Observe that for any $i\in I$,  $\#\partial K_i\le m-1$.
If a point $x$ lies in some $P_i$, then we call $x$ the {\em boundary point} for $K_i$.\\


\begin{definition}\label{refin}\cite{TYK}
Let  $ \mathcal K = \{K_i,i \in I =  \{1,\hdots,m \} \}$  and  $\mathcal L= \{L_j,j \in J =  \{1,\hdots,n \} \}$  be SIP systems of continua in $X$.\\
We say that $\mathcal L$ {\em is the refinement} of the system $\mathcal K$ if
\begin{itemize}
\item[(i)] for any $L \in \mathcal L$ there is $K \in \mathcal K$ such that $L \subset K$;
\item[(ii)]for any $K \in \mathcal K$ the union $|{\mathcal L}_K| $ of the sets of a subsystem $\mathcal L_K = \{ L \in \mathcal L : L \subset K \}$ is connected and contains $\partial K$.
\end{itemize}

\end{definition}

\begin{definition}\label{refsys}\cite{TYK}
Let $\{\mathcal K^n \} = \mathcal K^1, \mathcal K^2, \hdots $ be a sequence of SIP systems of continua. If for any $n$ the system $\mathcal K^{n+1}$ is a refinement of the system $\mathcal K^{n}$ and\\ for any nested sequence  $\{K^n_{i_n} \}_{n=1}^\infty$ such that $K^n_{i_n} \in \mathcal K^n$ and $K^n_{i_n} \supset K^{n+1}_{i_{n+1}}$,\\ the intersection $\bigcap\limits_{n=1}^\infty K^n_{i_n} = \{x\}$ is a singleton,\\ then we call $\{\mathcal K^n\}$ {\em a refinement sequence} for the set $K=|\mathcal K|$.
\end{definition}

\begin{proposition}\label{sip_qom}
Let  $\{\mathcal  K^n\}$  and $\{\tilde{\mathcal {K}}^n\}$ be   refinement sequences of SIP systems of continua for the sets $K$ and $\tilde{K}$, such that 
\begin{itemize}
\item[(i)] for any $n$, there is a bijection  $ \varphi:\mathcal{K}^n\to  \mathcal{\tilde K}^n$   such that $\varphi( K_i^n) =\tilde K_i^n $;
\item[(ii)] for any $K_i^n, K_j^n\in{\mathcal K^n}$, $K_i^n\cap K_j^n\neq\varnothing $ iff $ \tilde  K_i^n\cap \tilde K_j^n\neq\varnothing$;
\item[(iii)] for any $K_i^n\in\mathcal K^n,\ K_j^{n+1}\in\mathcal K^{n+1}$,\  $K_i^n\supset K_j^{n+1}$ iff $ \tilde  K_i^n\supset \tilde K_j^{n+1}$.
\end{itemize}
Then there exists a homeomorphism $f : K \to \tilde{K}$ such that for any $i$ and $n \in \mathbb{N} $ $f(K_{i}^n) = \tilde{K}_{i}^n$. 
\end{proposition}

\begin{proof}

For any $x\in K$ there is a  nested sequence  $K_{ i_n}^n$  such that $x = \bigcap\limits_{n=1}^\infty K_{ i_n}^n$. Put $f(x)=\tilde x= \bigcap\limits_{n=1}^\infty \tilde K_{ i_n}^n$. By   (iii), the intersection is a singleton, so $\tilde x$ is correctly defined. If $x=\bigcap\limits_{n=1}^\infty K_{ j_n}^n$ for some $ K_{ j_n}^n \neq K_{ i_n}^n$, then for any $n$, $K_{ i_n}^n\cap K_{ j_n}^n=\{x\}$   implies that $\tilde K_{ i_n}^n\cap \tilde K_{ j_n}^n=\{\tilde x\}$, so
the function $f$ is correctly defined. \\
Let $x = \bigcap\limits_{n=1}^\infty K_{ i_n}^n$ and
$  y = \bigcap\limits_{n=1}^\infty K_{ j_n}^n$ in $K$.\\
Clearly, $x \neq y $ iff there is $n \in \mathbb{N}$ such that 
$K_{ i_n}^n \cap K_{ j_n}^n = \varnothing$, which is equivalent to $\tilde{K}_{i_n}^n \cap \tilde{K}_{j_n}^n = \varnothing $ and  therefore  to $\tilde{x} \neq \tilde{y}$. This shows that the map $f$ is a bijection. \\

The relation $x\in K_i^n$ iff $\tilde x\in\tilde K_i^n$ implies that for any $K_i^n$, $f(K_i^n)=\tilde K_i^n$.

Given a point $x\in K$, define 
$V_n(x)= \bigcup\limits_{K_{i_k}^n \ni x_0} K_{i_k}^n $. By Definition \ref{refsys}, the family $\{V_n(x), n\in \mathbb{N}\}$ is a neighborhood base for $x$. The sets  
$f(V_n)=\bigcup\limits_{\tilde K_{i_k}^n \ni \tilde x} \tilde K_{i_k}^n=\tilde V_n$ form a
 neighborhood base for $\tilde x$. Therefore, $f$ is a homeomorphism. 
\end{proof}

\begin{theorem}\label{isom}
   If the systems ${\eS}=\{S_1,...,S_m\}$ and $\tilde {\eS}=\{\tilde S_1,...,\tilde S_m\}$ have isomorphic $P$-sprouts $\Gamma$ and $\tilde \Gamma$, then their attractors $K$ and $ \tilde K$ are isomorphic.
\end{theorem}

\begin{proof}

Consider two sequences of families of sets ${\EuScript K}^n=\{K_{\bf {i}}, {\bf {i}}\in I^{2^n}\}$ and $\tilde{\EuScript K}^n=\{\tilde K_{\bf {i}}, {\bf {i}}\in I^{2^n}\}$. Each of these families is the SIP system of continua. The families
 ${\EuScript K}^{n+1}, \tilde{\EuScript K}^{n+1}$ are the refinements of the systems  $ {\EuScript K}^n, \tilde{\EuScript K}^n$ respectively, so ${\EuScript K}^n, n\in\mathbb{N}$ and $\tilde{\EuScript K}^n, n\in\mathbb{N}$ are the refinement sequences of SIP systems of continua for the sets $K$ and $\tilde K$. The intersection graphs for the  systems  $ {\EuScript K}^n, \tilde{\EuScript K}^n$ are the isomorphic sprouts $\Gamma^{2^n}$ and $\tilde \Gamma^{2^n}$. By Proposition \ref{sip_qom}, there is a homeomorphism $f : K \to \tilde{K}$ such that for any   $n \in \mathbb{N} $ and $K_{\bf { i}}\in {\EuScript K}^n$ $f(K_{\bf {i}}) = \tilde{K}_{\bf {i}}\in \tilde{\EuScript K}^n$. 
 
 Now, if we take some $j\in I$ and $K_{\bf {i}}\in {\EuScript K}^n$, then we can represent $S_j(K_{\bf {i}})$ as a finite union of copies $K_{j\bf {i}\bf {k}}\in  {\EuScript K}^{n+1}$, to obtain $f(S_j(K_{\bf {i}}))= \tilde S_j(\tilde K_{\bf {i}})$. This shows that for any $x\in K$, $f(S_j(x))=\tilde S_j(f(x))$.
\end{proof}

% \newpage
\section{The ramification points and the main tree of a self-similar dendrite.}
 
If $K$ is a dendrite, then {\em the order ${\mathop{\rm Ord} \nolimits}(x,K)$ of a point} $x\in K$ is equal to the number of connected components of $K\setminus\{x\}$.  \cite[Ch.6,Th. 6]{Kur}. 
     
The points of order 1 are called the {\em end points} of $K$, the points of  order $\ge$ 2 are  the {\em cut points} of $K$,  the points of order $\ge$ 3 are the {\em ramification points} of $K$.


 \begin{definition}\label{mtree}\cite{STV}
Let the attractor $K$ of the  system ${\eS}=\{S_1,...,S_m\}$ be a self-similar dendrite with a finite self-similar boundary $\partial K$. 
The minimal subdendrite $\hga{\subset} K$ that contains $\partial K$ is called the {\em main tree} of $K$.
\end{definition}
Since  $\partial K=\{p_1,...,p_n\}$ is finite and for any two points $x,y\in K$ there is a unique subarc $\gamma(x,y){\subset} K$,
the main tree $\hga$ is a finite union of all subarcs $\gamma(p_i,p_j){\subset} K$. Thus, $\hga$ is
a finite topological tree. As was proved in
\cite[Theorem 15] {STV}, the arcs $\gamma(p_i,p_j)$ are the components of the attractor of some multizipper, which is a graph-directed
system that has the SIP. The arcs $\gamma(p_i,p_j)$ are called {\em the main subarcs} of the dendrite $K$.

We consider the set $G_{\eS}(\hga)$ that is the union of all images $S_{\bf j}(\hga)$ of the main tree, and its subset
$G_{\eS}(C)=\bigcup\limits_{{\bf j}\in I^*} S_{\bf j}(C)$, that is the set of all images of critical points $c\in C$.  It is essential because  
\begin{equation}\label{uadr}
 G_{\eS}(C)=\{x\in K: \#\pi^{-1}(x)>1\}.  
\end{equation}



\begin{proposition}\label{ord_cut}
Let the attractor $K$ of the system ${\eS}$ be a self-similar dendrite with finite boundary ${\partial} K$ and $\Gamma$ be its $P$-sprout.\\
1) %If $\Gamma({\eS})$ is regular, then 
For any $x\in K\setminus G_{\eS}(C)$, ${\mathop{\rm Ord} \nolimits}(x,K)\le\#P$.\\
2)  For any $x\in K\setminus G_{\eS}(C)$, there is ${\bf {i}}\in {I^*}$ such that ${\mathop{\rm Ord} \nolimits}(x,K)={\mathop{\rm Ord} \nolimits}(x,S_{\bf {i}}(\hga))$.\\
3) The set of cut points $CP(K)$ is contained in the set $ G_{\eS}(\hat \gamma)$.\\ 
4) Each  point $x\in K\setminus G_{{\eS}}(C)$ has a unique  address. If $x$ is a ramification point, this address is preperiodic.\\
5) $\hga{\subset} \bigcup\limits_{i=1}^m S_i(\hat \gamma)$.
\end{proposition}

\begin{proof} 1) Let  $x \in K\setminus G_{\eS}(C)$.  Let $Q_1,... Q_n$ be the connected components of $K \setminus \{x\}$. 
For each $k=1,...,n$ take  $x_k \in Q_k\setminus\{x\}$. There is ${\bf {i}}=i_1...i_n\in I^*$ such that
$x \in K_{\bf {i}}$ and
each $x_k\in K\setminus K_{\bf {i}}$. For each $k$ there is a unique arc $\gamma_k=\gamma(x,x_k){\subset} K$.


Note that $\gamma_k\setminus\{x\}{\subset} Q_k$ and
$\gamma_k\cap {\partial} K_{\bf {i}}=\{y_k\}$, where $y_k\neq x$. For different $k,l$,  
$\gamma(x,y_k)\cap \gamma(x,y_l)=\{x\}$, and all these arcs are  subarcs of
$S_{\bf {i}}(\hga)$. Therefore, $n\le\#P$.


To prove 2), take
$n={\mathop{\rm Ord} \nolimits}(x,K)$. Then ${\mathop{\rm Ord} \nolimits}(x,K)={\mathop{\rm Ord} \nolimits}(x,S_{\bf {i}}(\hga))$.

3) follows from 2). At the same time, $G_{\eS}(\hat \gamma){\supset} G_{\eS}({\partial} K)$; some points of the last set are the endpoints of $K$.

4) Since  $x\notin  G_{{\eS}}(C)$, the point $x$ has a unique address. 
By 2), there is a copy $K_{\bf {j}}$ such that  ${\mathop{\rm Ord} \nolimits}(x,K)={\mathop{\rm Ord} \nolimits}(x,S_{\bf {j}}(\hga))$. %For any boundary point $p\in{\partial} K$, the subarc $\gamma(x,p){\subset}\hga$ has a non-empty intersection 
%$\gamma(x,p)\cap{\partial} K_{\bf {j}}$, 
%therefore, 
%$ x\in \hga\cap K_{\bf {j}}{\subset} S_{\bf {j}}( \hga)$ and ${\mathop{\rm Ord} \nolimits}(S_{\bf {j}}^{-1}(x),\hga)\ge {\mathop{\rm Ord} \nolimits}(x,\hga)$. 
If $x$ is a ramification point,  its predecessor $ S_{\bf {j}}^{-1}(x)$ 
%$\in \hga\setminus G_{{\eS}}(C)$
is a ramification point of $\hga$ and lies in $\hga\setminus G_{{\eS}}(C)$.
Since the  number of  ramification points of $\hga$ is finite,  each ramification point $x\in K\setminus G_{{\eS}}({\partial} K)$ has a preperiodic address. 

5) Let $x\in\hga$ and $x\in S_i(K)$. If $x\notin {\partial} K_i$, there is a main subarc $\gamma(p_j,p_k)\ni x$. If any of these two points, say $p_j$, does not belong to ${\partial} K_i$, then $\gamma(x,p_j)\cap {\partial} K_i$ is a point  $\tilde p_j\in {\partial} K_i$. Thus, the arc $\gamma(p_j,p_k)\cap K_i$ is a  main subarc in  $K_i$. 
\end{proof}







To evaluate the order of the boundary points of $K$, we consider the sets of  addresses of these points.
For any $p\in {\partial} K$, the order ${\mathop{\rm Ord} \nolimits}(p,K)$ is equal to $\#\pi^{-1}(p)$ if the set $\pi^{-1}(p)$ is finite, and this order is infinite if $\pi^{-1}(p)$ is infinite. By Lemma \ref{ordpi},  $\pi^{-1}(p)=\hat\varphi( \Omega(p))$.

 \begin{definition}\label{cycles}\cite{BG} Let ${{\mathcal G}}$ be a finite digraph and $p, p_1$ be vertices in ${{\mathcal G}}$.
 
We write $p\prec p_1$ if $p$ is a predecessor of the vertex $p_1$ and $p_1$ is not a predecessor of the vertex $p$ in the digraph ${{\mathcal G}}$.
 
For a cycle $\sigma$ in ${{\mathcal G}}$ that does not contain a vertex $p$, we write $p\prec\sigma$, if there is a vertex $p_1$ in $\sigma $ such that $p \prec p_1$ and   $p_1\not\prec p$. We write $p\preceq\sigma$ if $p\prec\sigma$ or $p\in\sigma$.
    
For a pair  $\sigma_1,\sigma_2$ of disjoint cycles, we write $\sigma_1\prec \sigma_2$, if there are $p_1$ in $\sigma_1$ and $p_2$ in $\sigma_2$ such that $p_1\prec p_2$.
  
If a vertex $p_1$ in $\sigma_1$   is a predecessor of a vertex $p_2$ in $\sigma_2$ and $p_2$ is  a predecessor of the vertex $p_1$, then $p_1$ or $\sigma_1$ and  $p_2$ or $\sigma_2$ are {\em linked}.
  
If $\sigma_1\not\prec \sigma_2$ and  $\sigma_2\not\prec\sigma_1$, the  cycles $\sigma_1$ and $\sigma_2$ are {\em independent}. 
\end{definition}
   
Note that if a cycle $\sigma$ and a vertex $p$ outside $\sigma$ are linked, then there is a cycle $\sigma'$ that contains the vertex $p$ and some vertex $p'$ from the cycle $\sigma$, and therefore $\sigma$ and $\sigma'$ are linked. \\
If some cycles   $\sigma$ and $\sigma'$ are linked, then for any point $p$ in $\sigma$, the set of all walks  $\omega(p,...)$ that contain the points from $\sigma$ and $\sigma'$, is uncountable.
  
\begin{proposition}\label{numwalks}
Let ${{\mathcal G}}_P=(P,{{\mathcal E}}, \hat\varphi)$ be the index diagram of  $\Gamma({\eS})$ and $p\in P$.
\begin{itemize}
\item[1)] If all the cycles $\sigma_k$ in ${{\mathcal G}}_P$ for which $p\prec\sigma_k$, are independent, then all infinite walks $\omega(p,...)$ in  ${{\mathcal G}}_P$ are preperiodic and the set $\Omega(p)$ is finite for any $p\in P$. There is a uniform bound $M$ for $\#\Omega(p), p\in P $.
\item[2)] If there are cycles  $\sigma_1$ and $ \sigma_2$ in ${{\mathcal G}}_P$ such that $p\preceq \sigma_1$  and  $\sigma_1\prec \sigma_2$, then the set $\Omega(p)$ is infinite.
\item[3)] If $p\preceq\sigma_0$, $p\preceq\sigma_1$ and $\sigma_0,\sigma_1$ are linked, then $\Omega(p)$ is uncountable.
\item[4)] If $p\notin\sigma$ and  $p$ is linked to $\sigma$, then $\Omega(p)$ is uncountable.
\end{itemize}
\end{proposition}

\begin{figure}[htp]
\centering
\includegraphics[width=.9\linewidth]{pictures/infram_index.pdf}
\caption{A self-affine dendrite $K$ (left), its $P$-sprout (center) and the index diagram (right). The set of addresses of $p_2 \in \partial K$ in uncountable by  Proposition \ref{numwalks}/3. Each of the points $p_1, p_3 \in \partial K$ has a  unique address by   Proposition \ref{numwalks}/1. }
\label{fig:infram_index}
\end{figure}

\begin{proof} 
Each infinite walk $\omega=\omega(p_0,...)$ in ${{\mathcal G}}$ contains some cycle $\sigma$. If $\omega$ leaves $\sigma$ at some point $p$, then its subwalk $\omega'(p,...)$,  contains some cycle $\sigma'$. 
Then $\sigma\prec\sigma'$, or $\sigma$ and $\sigma'$ are linked. If several paths $\sigma_1,...\sigma_k$ are indepenent, then $\omega=\omega(p_0,\sigma_i)$ for some $i$.


If $p\prec  \sigma$ and there is no cycle $\sigma'$ such that $p\prec \sigma'\prec \sigma$ or $p\prec \sigma\prec \sigma'$, then each infinite walk $\omega(p,\sigma)$  is a composition of a path  $\omega_0(p,p')$ for some point $p'\in\sigma$ and an infinite walk $\omega_1(p',...)$    in $\sigma$. There is a finite number of possible choices of $\omega_0(p,p')$ for a given $p'$ and of  possible choices of $p'\in \sigma$. 
Applying the argument to each of the cycles $\sigma_k$, we obtain 1).



Let $p\preceq\sigma_1\prec\sigma_2$.  Then for each path $\omega(p,\sigma_2)$
 that has a non-empty intersection with $\sigma_1$, there are the vertices $p_1, p_2$ in $\sigma_1$ and a vertex $p_3$ in $\sigma_2$ such that $\omega(p,\sigma_2)$ is a composition of 4 paths: 1. $\omega(p,p_1)$;  2. a   path $\omega(p_1,p_2)$ in the cycle $\sigma_1$;  3. $\omega(p_2,p_3)$; and 4. an infinite path $\omega(p_3,...)$ in the cycle $\sigma_2$. There are infinitely many choices for $\omega(p_1,p_2)$ and $\omega(p_2,p_3)$, which proves 2).

If $\sigma_0$ and $\sigma_1$ are linked, take two vertices $q_i$ , $i=0,1$ that are contained in $\sigma_i$ and not  in $\sigma_{1-i}$.
For any infinite sequence $\alpha=i_1i_2...$ of digits 0 and 1 there are the paths $\omega(p,q_{i_1})$, $\omega(q_{i_k},q_{i_{k+1}}), k\in \mathbb{N}$, that are contained in $\sigma_{i_{k}}$ if $i_k=i_{k+1}$.

The composition of these paths is an infinite path $\omega_\alpha(p,...)$ in ${{\mathcal G}}$. For different values of $\alpha$ the paths $\omega_\alpha(p,...)$ are different, therefore, the set $\Omega(p)$ is uncountable.

The same argument applies to Case 4).
\end{proof}


Note that a regular sprout can be realized as a self-similar set in $\mathbb R^n$ only in Case 1).
But there is an example of a self-similar dendrite in the Hilbert space \cite{AK2025}, containing a boundary point $p$ with an infinite set $\Omega(p)$ and corresponding to an irregular sprout.

\begin{theorem}\label{inf_ram}
\qquad
\begin{itemize} 
\item[1)] If the index diagram of a sprout $\Gamma$ does not contain cyclic  vertices of outdegree $\ge 2$,
then the  orders of the points of $K$ are uniformly bounded.
\item[2)] If for some vertex $p \in  P $ there is a walk $\omega(p,..)$ such that all  vertices $p_i\in \omega(p,..)$ have  outdegree 1, then all  vertices $p_i\in \omega(p,..)$ have a unique preperiodic address.
\item[3)] Let $p$ be a common vertex of  two cycles  of the index diagram ${{\mathcal G}}_P$. Then, for any vertex $p'\prec p$, the set of infinite paths is uncountable.
\end{itemize}
\end{theorem}
  
\begin{proof} 
1) By Proposition \ref{numwalks},  for any $p\in {\partial} K$, ${\mathop{\rm Ord} \nolimits}(p,K)\le M$.
For any $x\in {\mathcal C}$, ${\mathop{\rm Ord} \nolimits}(x,K)\le m M$.
For any $x\in G_{\eS}({\partial} K)$, there is a multiindex ${\bf {i}}\in  I^*$ of maximum length such that  $x\in K_{\bf {i}}\setminus {\partial} K_{\bf {i}}$. Then $x=S_{\bf {i}}(y)$ for some 
$y\in{\mathcal C}$. Therefore, $\#\Omega(x)=\#\Omega(y) \le mM$. 

2) Since the outdegree of any vertex in $\omega(p,..)$ is 1, there is a cycle $\sigma$ such that $\omega(p,..)=\omega(p,\sigma)$, so $\#\Omega(p)=1$.

3) By Proposition \ref{numwalks},  the set $\Omega(p)$ is uncountable. 
Therefore, for any $p'\prec p$  the set  $\Omega(p')$ is uncountable.
\end{proof} 

\section{Finding the order of a point $x\in\hga$ with respect to the main tree.}

In the following, we  assume that $\Gamma=\Gamma(W,B,E,\varphi)$ is a {\em regular} $P$-sprout   with the sets of white vertices $W=\{w_1,..,w_m\}$, of indices $I=\{1,..,m\}$,  of black vertices $b\in B$ and of boundary points $p\in P$.\\






For each subset $P'$ of the set $P$, there is the smallest subtree  $\gamma(P') \subset \hga$ that contains the subset $P'$. 
The set of all the  subtrees  $\gamma(P')\subset \hga$ will be denoted by $\mathcal{P}(\gamma)$ or simply by $\mathcal{P}$.\\
The smallest subtree of the graph $\Gamma=\Gamma(P)$ that contains the subset $P'$ will be denoted by $\Gamma(P')$.
A subset $P'\IN P$ is {\em full} if $\gamma(P') \cap P=P'$, which is equivalent to $\Gamma(P') \cap P=P'$. By default, we consider only full subsets $P'\IN P$.


{\bf Definition of the mapping $\phi_i$.} Given $w_i\in  W$, we define the mapping $\phi_i:P\to P$ as follows.
For any   $p\in P$  there is a unique path $\beta(p,w_i)$  in $\Gamma$ with endpoints $p$ and $w_i$.
This path is an alternating sequence of  black and white vertices  $b_{i_0} w_{j_0}....b_{i_k} w_{j_k} $ 
 where $b_{i_0}=p$ and $w_{j_k}=w_i$,  such that each two consecutive elements are the endpoints 
 of some edge in $\Gamma$. Let the last edge of this path be $e=(b_{i_k} w_{j_k})$   and  let $\varphi(e)$ be the label on $e$. Then  put $\phi_i(p):=\varphi(e)$.\\
 
 There may be a degenerate case, where some $p_k\in P$ is incident to $w_i$, then $p_k\in K_i$ and $\phi_i(p_k)=S_i^{-1}(p_k)$.\\



We use the notation $\#_c(A)$ for the cardinality of the set of connected components of a set  $A$.\\

\begin{lemma}\label{degwi}
For any $i\in I$,  $\#\phi_i(P)= \deg(w_i, \Gamma)$ and
$$\#_c(K\setminus K_i)=\#_c(\hga\setminus K_i)=\#\phi_i(P) -\#(K_i\cap \partial K). 
$$
\end{lemma}
\begin{proof}
By its definition, $\phi_i(P)$ is the set of labels on the edges, incident to $w_i$. Since $\Gamma$ is regular, 
$\Gamma\setminus \{w_i\}$ is a disjoint union of components $C_k$, each containing one of the edges $e_k$ and a nonempty subset $P_{e_k}\IN P$ such that $P=\bigsqcup P_{e_k}$.

 In the  degenerate
case, a component $C_k$ contains only an edge $e_k=(p_k,w_i)$ and the boundary point $p_k$, and $P_{e_k}=\{p_k\}$. In that case $p_k\notin K\setminus K_i$ and $p_k\notin \hga\setminus K_i$.

If $C_k$ is non-degenerate $p_k\notin P_{e_k}$, there is a component of $\hga\setminus K_i$ and a component of   $K\setminus K_i$ that contain $P_{e_k}$ and $S_i(\phi_i(p))$.
    
\end{proof}
 

\begin{figure}[htp]\label{Fy1} 
\centering
\includegraphics[width=.80\textwidth]{pictures/graph_phi.pdf}
\caption{}
\end{figure}

For the transformation $\phi_i:P\to P$,  denote the set $\phi_i(P)$ by $P_i$,  the subgraph $\Gamma(P_i)$ by $\Gamma_i$ and  the subtree $\gamma(P_i)$ by $\hga_i$.\\



Notice that $\Gamma_i$ is the minimal subgraph of $\Gamma$ that contains $P_i$ and $\hga_i$ is a minimal topological subtree of $K$ that contains $P_i$.\\

{\small For the sprout $\Gamma$ in Figure \ref{Fy1},  $\phi_1 = \begin{pmatrix}
1 & 2 & 3 & 4  \\
2 & 4 & 4 & 1
\end{pmatrix} $ and $\phi_1(P) = \{p_1,p_2,p_4\}$. We show $\Gamma_1$ in the top-right, $\hga_1$ in  the bottom-left, and  
$S_1(\hga_1) = K_1 \cap \hat{\gamma}$ in  the bottom-right of the Figure.}

The definition of   map $\phi_i$ has an equivalent formulation in terms
of subarcs of the main tree $\hga$, which easily extends to the maps $\phi_{\bf j}$ for any ${\bf j}\in I^*$.


\begin{lemma}
Let $\gamma(p, K_i)$  be the
 minimal arc connecting $p$ and $S_i(K)$ in  $K$.  
 Then its endpoint in $K_i$,  i.e. $\gamma(p, K_i)\cap\partial K_i$, is $S_i(\phi_i(p))$. \hfill $\square$
\end{lemma}


 In other words, $\phi_i(p)=S_i^{-1}(\gamma(p, K_i)\cap\partial K_i)$. 
 
 Also, for any full $P'\subset P$, 
$S_i(\phi_i(P')) \subset\gamma(P')\cap \partial K_i$. In particular,

 $S_i(P_i)=\hga\cap \partial K_i=S_i(\hga_i) \cap \partial K_i$.\\
 
 
 
 

\begin{proposition}\label{fybj}
 Let ${\bf j}=j_1...j_k\in I^k$ and put $\phi_{\bf {j}}=\phi_{j_k}\cdot...\cdot\phi_{j_1}$.\\
Let $\gamma(p, K_{\bf j})$  be the
 minimal arc connecting $p$ and $S_{\bf j}(K)$ in  $K$.  \\
 Then its endpoint in $K_{\bf j}$  is $\gamma(p, K_{\bf j})\cap\partial K_{\bf j}=S_{\bf j}(\phi_{\bf j}(p))$.
 
 \end{proposition}
\begin{proof}
Let $i,j\in \{1,....,m\}$. Consider the minimal arc $\gamma(p,K_{ij})$. 
This arc is a union of minimal arcs $\gamma(p, K_i)$ and $S_i(\gamma(\phi_i(p), K_{j})) $.  
These arcs have a common endpoint $S_i(\phi_i(p))$. 
Since 
$\gamma(\phi_i(p), K_{j})\cap K_j=\{S_j(\phi_j\cdot \phi_i(p))\}$, 
the endpoints of $\gamma(p,K_{ij})$  are $p$ and $S_{ij}(\phi_j\cdot\phi_i(p))$. 
Therefore,  $\phi_{ij}=\phi_j\cdot\phi_i$.\\
By induction, we ensure that for each ${\bf {j}}=j_1...j_k$,  $\phi_{\bf {j}}=\phi_{j_k}\cdot...\cdot\phi_{j_1}$. 
For a minimal arc $\gamma(p,K_{\bf j})$ 
we obtain $\gamma(p,K_{\bf j})\cap K_{\bf j}=S_{\bf j}(\phi_{\bf j}(p))$. \end{proof}

 

{\bf Notation: $P_{\bf j},\Gamma_{\bf j}$ and $\hga_{\bf j}$.} For any full $P'\subset P$ and ${\bf{j}}\in I^*$,\\
$S_\bj(\phi_\bj(P'))=\gamma(P')\cap \partial K_\bj$, therefore $S_{\bf{j}}(P_{\bf{j}})=\hga \cap \partial K_{\bf{j}}$.\\
Hence, we denote the set $\phi_{\bf{j}}(P)$ by $P_{\bf{j}}$,  the subgraph $\Gamma(P_{\bf{j}})$ by $\Gamma_{\bf{j}}$, and  the subtree $ \gamma(P_{\bf{j}})$ by $\hga_{\bf{j}}$.\\
As  follows from Proposition \ref{fybj},
$P_{\bf{ij}}=\phi_{\bf{ij}}(P)=\phi_{\bf{j}}(P_{\bf{i}})$.\\

 {\bf The mapping $\phi^*_i$}.  The 
 equality $\gamma(p,K_{\bf j})\cap K_{\bf j}=S_{\bf j}(\phi_{\bf j}(p))$ leads us to introduce a  mapping
$\phi^*_i:\mathcal{P}(\gamma)\to
\mathcal{P}(\gamma)$ defined by the formula\\
$\phi^*_i(\gamma(P'))=\gamma(\phi_i(P'))$ for each  full $P'\subset P$.\\
Rewriting its right-hand side   gives
$\phi^*_i( \gamma(P'))=S_i^{-1}( \gamma(P')\cap K_i)$. Proceeding  to multiindices, we obtain that for any   $\bf {j}$,
$S_{\bf {j}}(\phi_{\bf {j}}^*(\hga))=\hga\cap K_{\bf {j}}$, or, in short,   $\phi^*_{\bf {j}}(\hga)=\hat \gamma_{\bf {j}}$.



   If $x, y \in \hat{\gamma}$ and $y =S_{\bf{j}}(x), \bf{j} \in $ $I^*$, then 
   \begin{equation}\label{eq:ord}
{\mathop{\rm Ord} \nolimits}(y,\hga) \geq {\mathop{\rm Ord} \nolimits}(y,K_{\bf{j}} \cap \hga) =  {\mathop{\rm Ord} \nolimits}(x,\phi^*_{\bf{j}}(\hga)).
\end{equation}
Therefore, for any ${\bf i,j}\in I^*$,
      $\#\phi_{\bf{i}}(P) \geq \#\phi_{\bf{ij}}(P) $.


\begin{proposition}\label{ordim}
 % Let $\Gamma$ be a regular $P$-sprout, 
 
  For $b\in B$,  $j\in I$, ${\bf{i}}\in I^*$, the following properties are fulfilled 
  
  1) If $\#\phi_{{\bf{i}}j}
    (P)=1$, then $w_{j} \notin \Gamma_{\bf{i}}$,

   
  2) If $\#\phi_{{\bf{i}}j}
    (P)=k>1$, then $\deg(w_{j},\Gamma_{\bf{i}}) = k$,
    
     
    3)  If $\deg(b, \Gamma) = k>1$, then $\Ord(b,\hga) \geq k$, 
    
4) If $\deg(b, \Gamma_{\bf{i}}) =k $, then $\Ord(y,\hga) \ge  k$, where $y = S_{\bf{i}}(b)$.
\end{proposition}
\begin{proof}
  
1) If $\#\phi_{{\bf{i}}j}(P) = 1$, then $\deg(w_j, \Gamma_{\bf{i}}) = 1$ therefore $w_{j} \notin\Gamma_{\bf{i}}$.   

 2) Since $\phi_{{\bf{i}}j}(P) = \phi_{j}( \phi_{\bf{i}}(P))$,   $\deg(w_{j}, \Gamma_{\bf{i}}) = \#\phi_{j}(\phi_{\bf{i}}(P))$.


3) If $\deg(b, \Gamma) = k$  there exist  subarcs  $\gamma(p_1,b), \gamma(p_2,b),\dots \gamma(p_k,b)$ of $\hat{\gamma}$ such that for any $i\neq j$, \  $\gamma(p_i,b)\cap \gamma(p_j,b) = \{b\}$. It follows that $\Ord(b,\hga) \geq k$.


4)  If $\deg(b, \Gamma_{\bf{i}}) \geq k > 1$ there are subarcs $\gamma(p_i,b), i=1,...,k$ of $\phi^*_{\bf{i}}(\hat{\gamma})$ with the same property as in 3). It follows that $\Ord(b,\phi^*_{\bf{i}}(\hat{\gamma})) \ge k$. 
By formula \eqref{eq:ord}, 
$\Ord(S_{\bf{i}}(b),\hat{\gamma}) \ge\Ord(b,\phi^*_{\bf{i}}(\hat{\gamma})) \geq k$. Therefore, $\Ord(y,\hat{\gamma}) \ge  k$.
By formula \eqref{eq:ord}, 
$\Ord(S_{\bf{i}}(b),\hat{\gamma}) \ge  \Ord(b,\phi^*_{\bf{i}}(\hat{\gamma})) \geq k$. Therefore, ${\Ord} (y,\hat{\gamma}) \ge  k$.
\end{proof}


{\bf Definition and properties of     $\Nfi(\al)$. }\ For an address $\al=j_1\ldots j_k\ldots $, we define $\al|_n=j_1\ldots j_n $. Consider the sets $P_{\al|_n}=\phi_{\al|_n}(P)$.
For any $n$, $\#P_{\alpha|_n}\ge \#P_{\alpha|_{n+1}}\ge 1$, and hence
$\lim\limits_{n\to\infty}\#P_{\alpha|_n}=\min \#P_{\alpha|_n}$ is an integer $\ge 1$, which we denote by $\Nfi(\al)$.



  
\begin{theorem}\label{ord_b}
%Let $(\Gamma,\varphi)$ be a regular $P$-sprout.
For any $x\in K$, the set $A_x=\{\alpha\in\pi^{-1}(x):\Nfi(\alpha)  > 1\} $ is finite.  $A_x=\varnothing$ if and only if $x\notin\hat\gamma$.
\begin{itemize}
   
     \item [1)] If $\#A_x=1$ and $x\notin  \partial K$, then  ${\mathop{\rm Ord} \nolimits}(x,\hga) =  \Nfi(\alpha) $.
      \item [2)] If $\#A_x=1$ and $x\in  \partial K$, then  
     ${\mathop{\rm Ord} \nolimits}(x,\hga) =  \Nfi(\alpha)  - 1$.
     \item [3)] 
If $\#A_x>1$, then 
    \begin{equation}
      {\Ord}(x,\hga)= \sum\limits_{\alpha\in A_x} (\Nfi(\alpha)  - 1). 
    \end{equation}
     \end{itemize}
\end{theorem}
 
\begin{proof}
1) If the point $x \in \hga \setminus\partial K$ and $ A_x =\{\alpha\}=$ $j_1j_2\hdots$, then $x \in \hga \setminus G_{{\eS}}(\partial K)$  and $\Nfi(\alpha)=k > 1$. There is  $m$ such that for the initial subword $\bf{j} =$ $ j_1j_2\hdots j_m$ of $\alpha$,   $K_\bj\cap\partial K=\varnothing$ and $\#\phi_{\bf{j}}(P) = k$. 


The same holds for any $\bf{j}'$ such that $ \bf{j} \sqsubset \bf{j}' \sqsubset \alpha$. Therefore, $K_{\bf{j}'} \cap \hga$  is a subtree of $\hga$ that has $k$ endpoints and a unique  point $x$. Therefore,  ${\mathop{\rm Ord} \nolimits}(x,\hat{\gamma}) = \Nfi(\alpha) = k$.






 An example of this situation is shown in Fig. \ref{fig:case3ex}.
\begin{figure}[htp]
    \centering
    \includegraphics[width=\linewidth]{pictures/case3ex.pdf}
    \caption{ Attractor $K$, its $P$-sprout  and the index diagram. The boundary point $p_3$ has two addresses. Its order in $\hat{\gamma}$ is equal to 2. }
     
    \label{fig:case3ex}
\end{figure}

2) In the case where $x \in \partial K$, $ A_x =\{\alpha\}=j_1j_2...$, and $\Nfi(\alpha)=k > 1$, there is  $m$ such that for the initial subword $\bf{j} =$ $ j_1j_2\hdots j_m$ of length $m$,   $K_{\bf{j}}\cap \partial K=\{x\}$ and $\#\phi_{\bf{j}}(P) = k$.

These two relations hold for any $\bf{j}'$ such that $ \bf{j} \sqsubset \bf{j}' \sqsubset \alpha$. Then $K_{\bf{j}'} \cap \hat \gamma $  is a subtree of $\hga$ that has $k-1$ endpoints and a unique cut  point $x$, because $x \in \phi_{\bf{j'}}(P) $. Therefore,  ${\mathop{\rm Ord} \nolimits}(x,\hat{\gamma}) = \Nfi(\alpha)-1 = k-1$.





% An example of this situation is shown in Fig. \ref{fig:hex3}.

% \begin{figure}[htp]
%     \centering
%     \includegraphics[width=\linewidth]{pictures/hex3.pdf}
%     \caption{
%         Applying the result to an example in Fig.\ref{fig:hex_polygon2}. 
%         $\pi^{-1}(x)=\{\overline{3}\}$,\; 
%         $\{y_1,y_2,y_3\}\subset\partial K$,\;
%         $\{x_1^1,x_2^1,x_3^1\}=\partial K_3$,\;
%         $\{x_1^2,x_2^2,x_3^2\}=\partial K_{33}$.}
%     \label{fig:hex3}
% \end{figure}



3) Let $\alpha^1,...,\alpha^n$ be the addresses of the point  $x$, such that for any   $i \in \mathbb{N}$, $\Nfi(\alpha^i) > 1$.
For a sufficiently large $m$,  a set $\{\bf{j}^1,\bf{j}^2...,\bf{j}^n\}$ of initial subwords of  length  $m$ for the  addresses $\alpha^1,...,\alpha^n$ satisfies the following conditions.\\ a) All these multiindices are different.\\ b) For each $\bf{j}^k$, $K_{\bf{j}^k} \cap \partial K = \{x\}$ if $x \in \partial K$ and $K_{\bf{j}^k} \cap \partial K = \{\varnothing\}$ if $x \in \hat{\gamma}\setminus \partial K$.\\ c) Each copy $K_{\bf{j}^k}$  does not contain any boundary or ramification point of  $\hga$ except $x$. 

Under these conditions, $\#\phi_{\bf{j}^k}(P) = \Nfi(\alpha^k)$.
Then for any $\bf{j}'^k $ such that $\bf{j}^k \sqsubset \bf{j}'^k \sqsubset \alpha_k$,\ \ 
$\#S_{\bf{j}'^k} (P) \cap \hat \gamma = {\mathop{\rm Ord} \nolimits}(x, K_{\bf{j}'^k }  \cap \hat \gamma) + 1 $.

Applying the above argument to each of the addresses of the point   $x$, we obtain the inequality 
 $\sum\limits_{k=1}^n {\mathop{\rm Ord} \nolimits}(x, K_{{\bf{j}^k}} 
\cap \hat \gamma) \le  {\mathop{\rm Ord} \nolimits}(x, \hat \gamma) \le \#P   $. Therefore, the set $A_x$ is finite and  ${\mathop{\rm Ord} \nolimits}(x,\hga) = \sum\limits_{\alpha \in A_x} (\Nfi(\alpha)  - 1)$.\\
If for some address $\alpha^i$ there is a multiindex $\bf{i}$ $\sqsubset \alpha^i$ such that $\#\phi_{\bf{i}}(P) = 1$, then $K_{\bf{i}} \cap \hat \gamma = \{x\}$. Consequently, ${\mathop{\rm Ord} \nolimits}(x, K_{\bf{i}} \cap \hga) = 0 $.\\
\end{proof}


%An example of Cases 2 and 3 is shown in Fig. \ref{fig:hex_polygon2}
\begin{figure}[htp]
\centering
\includegraphics[width=\textwidth]{pictures/hex_polygon2.pdf}
\caption{ An attractor $K$ with its main tree $\hat \gamma$, its $P$-sprout, and its index diagram. Each boundary point has a  unique address. ${\mathop{\rm Ord} \nolimits}(p_2, \hat \gamma) = \Nfi(2\overline{3}) -1= 2, {\mathop{\rm Ord} \nolimits}(p_3, \hat \gamma) = \Nfi(\overline{3}) -1= 3$. All the other points are endpoints. There is ramification point $a$, such that $\pi^{-1}(a) = 4\overline{3},\; {\mathop{\rm Ord} \nolimits}(a, \hat \gamma) = \Nfi(4\overline{3}) = 3 $ }
\label{fig:hex_polygon2}
\end{figure}


 \begin{corollary}\label{ram_adr}
        Let  $\alpha$ be the address of a point $x\in K$ and $\Nfi(\alpha) \ge 3$.\\ If $\pi^{-1}(x)=\{\alpha\}$, $x\in\partial K$ and $\Nfi(\alpha) = 3$, then ${\rm Ord}(x,\hga)=2$.\\ Otherwise,  ${\mathop{\rm Ord} \nolimits}(x,\hga) \ge 3$.
\end{corollary}

\begin{proposition}\label{ram_p}
       % Let $\Gamma$ be a regular $P$-sprout and  for some copy $K_{\bf{i}}$, where 
       If  for some $\bf{i} $, $\#\phi_{\bf{i}}(P) \ge 3$, then at least  one of the following statements holds:\\
1) $K_{\bf{i}}\cap \partial K$  is nonempty and contains a point  $x$ such that  ${\mathop{\rm Ord} \nolimits}(x,\hga) \geq 2$;\\
2) there is a point  $x \in K_{\bf{i}} \cap \hat{\gamma}$ such that  ${\mathop{\rm Ord} \nolimits}(x,\hga) \geq 3$.
\end{proposition}

\begin{proof}
If there is $x\in K_\bi\cap \partial K$, then there are $ \{a,b\} \IN (\partial K_\bi \cap \hga)\setminus\{x\}$. Consider the subarcs $\ga(a,x), \ga(b,x)$ of $\hga$. If $\ga(a,x)\cap \ga(b,x)=\{x\}$, ${\Ord}(x,  \hga) \ge 2$, and 1) follows. \\
Otherwise, there is $y$ such that $\gamma(a,x)\cap \gamma(b,x)=\gamma(y,x)$, and
${\mathop{\rm Ord} \nolimits}(y,  \hga) \geq 3$.\\
If $K_{\bf{i}}\cap \partial K=\varnothing$, take $ \{a,b,c\} \subset \partial K_{\bf{i}} \cap \hat{\gamma}$. 
 Consider the subarcs $\gamma(a,b)$, $\gamma(b,c)$, $\gamma(a,c)$ of $\hat{\gamma}$. The  intersection  $\gamma(a,b) \cap \gamma(b,c) \cap \gamma(a,c)$ is a singleton $\{x\}$.
 If $x \notin \{a,b,c\}$, ${\mathop{\rm Ord} \nolimits}(x, \gamma(a,b,c)) = 3$.\\
If $x \in \{a,b,c\}$, say $x=b$, then $\gamma(a,b) \cap \gamma(b,c) = \{b\}$. Since there is a subarc $\gamma'\subset\hga\setminus{K_{\bf{i}}}$  ending at  $b$, ${\mathop{\rm Ord} \nolimits}(x, \hga) \ge {\mathop{\rm Ord} \nolimits}(x,  (\gamma'\cup \gamma(a,b) \cup \gamma(b,c)) ) =3$.
\end{proof}
Proposition \ref{ram_p}  shows how to find  copies that contain the ramification points and thereby find the addresses of these points. In case the ramification points have multiple addresses, they are the images of boundary points. The addres\-ses of  boundary points are defined using the index diagram according to Pro\-po\-sition \ref{infwalks}. A tool for finding all the needed addresses will be a directed  graph $G_T$, which we define  below.\\



{\bf The graph $G_T$}. By analogy to the index diagram, we construct a labeled digraph $G_T=\{ V_T, E_T, \psi_T\}$,  such that each proper walk in this graph defines an address of a ramification point of $\hga$ or of a  boundary cut point in  $\hga$.\\


 The set of vertices  $V_T$ consists of three disjoint parts: $V_Q$,  
 $V_B$, and $V_P$. \\
 $V_Q$ is the set of all subsets
$Q\IN P$  such that  $\#Q\ge 3$ and for some $\bi\in I^*$, $\phi_\bi(P)=Q$; \\
$V_B=\{b\in B: \deg(b, \Gamma) \ge  3\}$;\\
$V_P=\{p\in P: \deg(p, \Gamma) \ge  2\}$. 
\\

The set of edges $E_T$ also consists of 3 disjoint parts, $E_Q$, $E_B$, and $E_P$.\\
The set $E_Q$ is supplied with a labeling function $\psi:E_Q\to I$.

A pair $Q,Q'\in V_Q$ defines  an edge $\be=(Q,Q')\in E_Q$ with a label $\psi(\be)=i$ exactly when there is $i\in I$ such that $\phi_i(Q)=Q'$. 

In the same way as in  the definition \ref{walks}, each walk $\be_1...\be_k$, $\be_i\in E_Q$,   has labels $l_j=\psi(\be_j),j=1,...,k$ and travels along a sequence  of vertices
 $Q_0,...,Q_k$ such that for any  $j=1,...,k$, 
 $\phi_{l_j}(Q_{j-1})=Q_j$. Consequently,  $\phi_{l_1....l_k}(Q_0)=Q_k$.
 
 In its turn, each infinite walk $\be_1...\be_k...$ in $E_Q$ that starts from $Q_0=P$ defines an address ${\beta }=l_1l_2...$ such that $\Nfi(\beta)\ge 3$.

 The sets of edges $E_B$ and $E_P$  are defined in the following way.
 
A pair $(Q,b)$, where $Q\in V_Q$ and $b\in V_B$, defines an edge $\be=(Q,b)\in E_B$ exactly when $b\notin Q$ and $\deg(b,G_Q)\ge 3$.

A pair $(Q,p)$, where $Q\in V_Q$ and $p\in V_P$, defines an edge $\be=(Q,p)\in E_P$ exactly when $p\in Q$ and $\deg(p,\Ga_Q)\ge 2$.\\
The edges $\be\in E_B\cup E_P$ have no labels;  they are the terminal edges  of finite paths in $G_T$.\\

A walk $\omega=(\be_1...\be_k...)$ is {\em proper}, if it starts from the vertex $P$, and cannot be extended to a larger walk.
In the following, we will consider only proper walks in $E_T$.

\begin{lemma}\label{ram_b}
    If for some vertex $b \in B$, $\deg(b, \Gamma) \geq 3$ then $b$ is a ramification point of $\hat{\gamma} $.
\end{lemma}
\begin{proof}
    By Proposition \ref{ordim}(3) ${\mathop{\rm Ord} \nolimits}(b, \hga) \ge 3 $. 
\end{proof}
The set of vertices $V_R = \{b \in V_B \cup V_P,\; \deg(b, \Gamma)\geq 3 \}$ will correspond to  the ramification points of $\gamma$ which are the intersection points of at least three copies or the boundary points of the critical set.

 \begin{proposition}\label{fin_cyc}
    Each vertex $Q \in V_Q$ of the graph $G_T$ has an outgoing edge. If $Q$ lies in some cycle, this edge is unique.
 \end{proposition}
 \begin{proof}
  If $Q \in V_Q$, then $\#Q = k \geq 3$. Then two cases are possible:\\
  (i) For any $p \in Q$,\; $\deg(p, \Gamma_{Q}) =1$. In this case, the graph $\Gamma_{Q}$  necessarily has a vertex $v\in W\cup B$ such that $\deg(v, \Gamma_{Q}) = k \geq 3$.\\ 
  If $v\in W$, then $\#\phi_v(Q)=k$. Hence, there is an edge $e=(Q,\phi_v(Q))\in E_Q$.\\
If $v\in B$, then 
  there is an edge $e=(Q,b) \in E_B$ in the graph $G_T$.\\
  (ii) If there is a vertex $p \in Q$ such that $\deg(p, \Gamma_{Q})\ge 2$, then there is an edge $e =(Q, p) \in E_P$ on the graph $G_T$.

  If $Q\in V_Q$ lies in some cycle, then for any vertex  $Q_j$ in that cycle $\#Q =\#Q_{j} = k \geq 3$. Let $Q_{j}=\phi_{j}(Q)$, then  $\deg(w_{j}, \Gamma_{Q}) = k $ and for any $p \in Q\;$ $\deg(p, Q) = 1$, so  $w_{j}$ is a unique vertex of order $\ge 3$ in $\Gamma_{Q}$. Therefore, $Q$ has a unique outgoing edge.
\end{proof}
     

 




In the transformation graph $G_T$, there are 3 types of proper walks. They produce the addresses of the ramification points of $\hga$.\\


{\bf  1.}   $\Omega_Q$  denotes the set of  proper walks whose edges belong only to $E_Q$.

\begin{lemma} \label{inf_walks}
Each walk $\omega=(\be_1,\be_2,...)\in\Omega_Q$ is infinite and terminates in a cycle.
It defines a preperiodic address ${\bm\beta}=\beta_1\beta_2...$, where $\beta_k=\psi(\be_k)$ such that $\pi({\bm\beta})=x\in\hga$, $\Nfi({\bm\beta}) \ge 3$,
and one of the following conditions  holds:\\
1) $x\in \partial K$ and $x$ is a cut point of $\hat \gamma$;\\
2) $x\notin \partial K$ and $x$ is a ramification point of $\hat \gamma$.
\end{lemma}

 \begin{proof}\label{path_bk}
 As graph $G_T$ is finite and by Proposition \ref{fin_cyc} at least one edge comes from each vertex $Q \in V_Q$, there exist $k, l \in \mathbb{N}$ such that $\bm{e}_k = \bm{e}_{k+l}$. It follows that infinite walk $\omega$ terminates in a cycle and for any of the  edges of $\omega$, $\psi({\bm e_j})= {{\beta}}_j$. Then $\omega$ defines a preperiodic address $\bm{\beta}={\beta}_{1}{\beta}_{2} \hdots \overline{{\beta}_{k+1} \hdots {\beta}_{k+l}}$. 

 The point $x = \pi({\bm\beta})$ belongs to  $\hat{\gamma}$ because for any finite subword ${\bm \beta}|_{l}$ of ${\bm{\beta}}$ of length $k$ $\#\phi_{\bm{\beta|_{k}}}(P)\geq 3$. There exist $p_1, p_2 \in \partial K$ such that for $j=1,2\;$
the minimal arc $(p_j, K_{\bm{\beta|_{k}}}) \in \hat{\gamma}$. Let $(p_j, K_{\bm {\beta}}) = \lim\limits_{k \to \infty} (p_j, K_{\bm {\beta|_{k}}})$. The arc $(p_j, K_{\bm {\beta}}) \in \hat{\gamma}$ and contains the point $x$ since ${\bm \beta}$ is an address of $x$. Therefore, $x \in \hat{\gamma}$. \\
  Since $\Nfi({ \bm \beta}) = k \geq 3$   by Corollary \ref{ram_adr}  a point $x \in \hat{\gamma}$ with address $\bm{\beta}$ is a boundary cut-point of $\hat{\gamma}$  or a non-boundary ramification point of $\hat{\gamma}$.
 \end{proof}
 

{\bf 2.} $\Omega_B$ denotes a set of proper finite paths $\omega_b$ in $E_Q$ that terminate with an edge $\be_b=(Q,b )\in E_B$. They  correspond to ramification points of $\hga$ that are images of  $b $. 

\begin{lemma}\label{path_bk}
    Let $\omega=(\be_1,...,\be_k,\be_b)$ be a path in $\Omega_B$. Let  $j_i=\psi(\be_i)$, $\bj=j_1...j_k$ and ${\bm e}_b = (Q_{i_k}, b_l)$.  
   Then  the point $x = S_\bj(b_l)$ is a ramification point of  $\hat \gamma$ and $\pi^{-1}(x) = {\bm j}\pi^{-1}(b_l)$.

\end{lemma}
 \begin{proof}
 If  $\deg (b_l, \Gamma_{Q_{i_k}}) = l \geq 3$, by Proposition \ref{ordim}(5)\\ ${\mathop{\rm Ord} \nolimits}(S_{\bm j}(b_l), \hat \gamma)\geq {\mathop{\rm Ord} \nolimits}(b_l, \phi^*_{\bm j}(\hat \gamma)) \geq l \geq 3$. Then the point $x = S_{\bm j}(b_l)$ is a ramification point of  $\hat \gamma$. Seeing that $b_l \notin Q_{i_k}$ we get  $x = S_{\bm j}(b_l) \notin \partial K_{{\bm j}}$. Сonsequently, $\pi^{-1}(x) = {\bm j}\pi^{-1}(b_l)$.
 \end{proof}



{\bf 3.} $\Omega_P$ denotes a set of proper finite paths $\omega_b$ in $E_Q$ that terminate with an edge $\be_p=(Q,p )\in E_P$. They  correspond to the boundary cut points or ramification points of $\hat{\gamma}$ that belong to the critical set.

\begin{lemma}\label{path_pk}
    Let $\omega=(\be_1,...,\be_k,\be_p)$ be a path in $\Omega_P$. Let  $j_i=\psi(\be_i)$, $\bj=j_1...j_k$ and ${\bm e}_p = (Q_{i_k}, p_l)$.  
   Let $x = S_\bj(p_l)$. Then\\
   1) If $x\in \partial K$ then $x$ is a cut point of $\hat \gamma$.\\
2) If $x\notin \partial K$ then $x$ is a ramification point of $\hat \gamma$.

\end{lemma}


 \begin{proof}
 1) If  $\deg (p_l, \Gamma_{Q_{i_n}}) = l \geq 2$, there will be $l$ arcs $\gamma_1, \dots \gamma_l$. Each $\gamma_i \subset \hat{\gamma}$ connects $x$ and $y_i \in \partial K_{{\bj}}$. Therefore ${\mathop{\rm Ord} \nolimits}(x, \hat \gamma) \geq l \geq 2$.

 2)  If  $\deg (p_l, \Gamma_{Q_{i_n}}) = l \geq 2$, there will be $l$ arcs $\gamma_1, \dots \gamma_l$. Each $\gamma_i \subset \hat{\gamma}$ connects $x$ and $y_i \in \partial K_{\bj}$. Since $x= S_{\bj}(p_l)$,  $x\in \partial K_{\bj}$. It means that $x\notin \partial K$ and $x\in K_{\bj^{-} i_n}$, where $\bj^{-} = j_1\hdots j_{n-1}$ and ${\mathop{\rm Ord} \nolimits}(x, K_{\bj^{-} i_n} \cap \hat{\gamma}) > 0$. Therefore ${\mathop{\rm Ord} \nolimits}(x, \hat \gamma) \geq l+1 \geq 3$.

 \end{proof}




  
\begin{lemma}
       In the transformation graph $G_T$ \; $\#(\Omega_Q \cup \Omega_B \cup \Omega_P) \le M  -2$, where $M = \#P$.
\end{lemma}\label{fin_walks}
\begin{proof}
    By \ref{fin_cyc}, any two cycles $\sigma_1, \sigma_2$ in $G_T$ are independent. From  Proposition \ref{numwalks} it follows that the set $ \Omega_Q$ is finite. Since the graph $G_T$ is finite, the number of paths in this graph is finite. Therefore, $\Omega_B$ and $\Omega_P$ are also finite. If for some $w_i \in W$ $\#\phi_i(P) = k\geq 3$ then graph $\Gamma_i$  may contain at most $k-2$ white vertices with a degree greater than or equal to 3.
\end{proof}






Let $I_{\partial K}$ be the set of addresses of the boundary points and $I_{C}$ be the  set of addresses of the critical  points. If for some address $\alpha =\beta \alpha_i,\; \alpha_i \in I_{C}, \beta \in I^* $, $\Nfi(\alpha) \geq 3$ then $x= \pi(\alpha)$ is a point of $\hat{\gamma}$ with multiple addresses.

 If  $ \alpha \in I_{\partial K} $ and $\Nfi(\alpha) \geq 3$ then $x= \pi(\alpha)$ is a boundary cut point of $\hat{\gamma}$.\\

 If the address $\alpha$ does not correspond to a boundary point or a point with multiple addresses, then $x = \pi(\alpha)$ is a point with a single address $\alpha$.




\begin{figure}[htp]\label{ram_tree}
\centering
\includegraphics[width= \textwidth]{pictures/ram_tree.pdf}
\caption{ An attractor $K$ and its main tree $\hat \gamma$ on the left. Its $P$-sprout on the right. The points $a, b, c, d $ are the ramification points of the main tree. $\pi^{-1}(a) = \overline{34},\; \pi^{-1}(b) = 2\overline{34},\;   \pi^{-1}(c) = \overline{43}, \;\pi^{-1}(d) = 5\overline{43}$.}
\label{fig:ram_tree}
\end{figure}


\begin{figure}[htp]\label{fig:gt}
\centering
\includegraphics[width=.8 \textwidth]{pictures/gt.pdf}
\caption{The transformation graph for Figure \ref{fig:ram_tree}}

\end{figure}



\begin{theorem}\label{radr}
    Let $G_T$ be the transformation graph of the $P$-sprout $\Gamma$. Then, for each ramification point $x$ of the main tree $\hat{\gamma}$, which is not contained in $B\cup P$, at least one of the following conditions holds:

1) There is a walk $\omega \in \Omega_Q$ that defines a preperiodic address ${\bm\beta}=\beta_1\beta_2...$,  such that $\pi({\bm\beta})=x$. This walk consists of some path $(\be_1,...,\be_k)\in \Omega_Q$ and a cycle
$(\be_k,...,\be_n)\in \Omega_Q$.

    
2) There is a path $\omega=(\be_1,...,\be_k,\be_b)\in \Omega_B$ with an end point $b$, which defines a multiindex $\bj=j_1...j_k$, where  $j_i=\psi(\be_i)$, such that  
    $x = S_{\bj}(b)$ .
    
3) There is a path $\omega=(\be_1,...,\be_k,\be_p)\in \Omega_p$ with an endpoint $p\in P$, that defines a multiindex $\bj=j_1...j_k$, where  $j_i=\psi(\be_i)$, such that  
    $x = S_\bj(p)$.




If $x$ is a ramification point of $\hat{\gamma}$ with a unique address, then only  1) is valid.
If $x$ is a ramification point with several addresses, then any of  statements 1) -- 3) may be valid.
\end{theorem}



\begin{proof}  Let $x$ be a ramification point of $\hat\gamma$ and let $\bmb={\beta_1\beta_2....}$ be its address.

Since $x$ belongs to a unique copy $K_{\beta_1}$ of order 1,   there is a copy $K_{\bf j} \ni x$, where ${\bf j} ={\beta_1\hdots \beta_n}$, for which all vertices $Q_i$ on the path 
     $P \overset{\beta_1}{\to} Q_{1} \overset{\beta_2}{\to} Q_{2}  \dots  \overset{\beta_n}{\to} Q_{n}$ in $G_T$    are different. The graph $\Gamma_{Q_n}$ satisfies one of the three following conditions.\\

     1)  There is a vertex $w_{\beta_{n+1}} \in W$ such that $\deg(w_{\beta_{n+1}},\Gamma_{Q_n})\ge 3$ and $x \in K_{{\bf j}\beta_{n+1}}$. Then there is an edge $Q_n\overset{\beta_{n+1}}{\to}Q_{{n+1}}$.
     
      If $Q_{n+1}$ is different from all $Q_i,i=1,...,n$, we repeat the induction step for the copy $K_{\bj'}\ni x$ and $\bj'=\beta_1...\beta_n \beta_{n+1}$.

     If $Q_{n+1}=Q_i$ for some $i\le n$, then the vertices $Q_i,...,Q_n$ form a cycle.  By proposition \ref{fin_cyc}, each of the vertices $Q_k$ of this cycle has a unique outgoing edge in $G_T$, which is defined by a unique $\beta_{k+1}$.    Therefore,
     $Q_{n+1+k}=Q_{i+k}$  and  $\beta_{n+1+k}=\beta_ {i+k}$ for any $k \in\mathbb{N}$.

     Consequently, by the finiteness of the set $V_Q$, each infinite walk in $\Omega_Q$ is preperiodic, which proves 1).
    
    2) There is a vertex $b_k \in B$, $\deg(b_k, \Gamma_{Q_{n}}) \ge 3$. Then there is an edge $Q_{n} \to b_k$ in the transformation graph $G_T$.
By  lemma \ref{path_bk}, there is a ramification point $x = S_{{\bf j}}(b_k)$ such that ${\bf j}$ is the initial word of all addresses of $x$. Therefore, $\pi^{-1}(x) = {\bf j}\pi^{-1}(b_k)$, which proves 2). \\

 3) There is a vertex $p_k \in Q_{n}$, $\deg(p_k, \Gamma_{Q_{n}}) \ge 2$. Then there is an edge $Q_{n} \to p_k$ in the transformation graph $G_T$. By  lemma \ref{path_pk}, there is a ramification point $x = S_{{\bf j}}(p_k)$ such that $\pi^{-1}(x) \supseteq {\bf j}\pi^{-1}(b_k)$.\\

 


If for a point $x$ condition 2 or 3  is satisfied, $x$ is a point of the critical set or its image and therefore has more than a  unique address. Сondition 1 is necessarily fulfilled, since for the unique address $\alpha$ of $x\;$, $N_{\phi}(\alpha) \geq 3$ and therefore there exists a walk $\omega(P,\dots) \in \Omega_Q$.


\end{proof}

As follows from the above, to find the addresses of all ramification points of  $\hat\gamma$, it is necessary to find all the cycles in the transformation graph $G_T$.
 By proposition \ref{fin_cyc}, all these cycles are independent. Since each of them produces a ramification  point or a boundary cut point, their number is  not greater than $\#P-2$. For each of these cycles one should find the set of all paths in $E_Q$ connecting  $P$ and this cycle.  Thus, we obtain the set $\Omega_Q$.
To find the sets   $\Omega_B$ and $\Omega_P$, 
one should find the set of all the paths in $G_T$ 
that connect $P$ and $b \in V_B$ or $p \in V_P$, respectively. 
 






\section*{Acknowledgments}
The authors especially thank Dmitry Mekhontsev, whose package IFSTile \cite{IFStile} was the main  tool  for the construction  of self-similar and self-affine sets shown in the paper.\\
The study was carried out under the state contract of the Sobolev Institute of Mathematics (project FWNF-2026-0026).
\footnotesize
\begin{thebibliography}{00}

\bibitem{AST2024}
 K. Allabergenova, M. Samuel, A. Tetenov,
 {Intersections of the pieces of self-similar dendrites in the plane},
 {Chaos, Solitons \& Fractals}, \textbf{182} (2024), 114805.

\bibitem{AK2025}
 K. Allabergenova, M. Kadirova,
 {On self-similar dendrites in Hilbert space},
 \emph{Results in Nonlinear Analysis}, \textbf{8}:1 (2025), 184--192.

\bibitem{Bandt2024}
 Ch. Bandt,
 {Elementary fractal geometry. 4. Automata-generated topological spaces},
 \emph{Communications in Mathematics} \textbf{33}:2:4 (2025), 1--33.

\bibitem{SSS2}
 Ch. Bandt and K. Keller,
 {Self--Similar Sets 2. A Simple Approach to the Topological Structure of Fractals},
 \emph{Mathematische Nachrichten} \textbf{154}:1 (1991), 27--39.

 \bibitem{SSS6} 
Ch. Bandt and J. Stahnke. { Self-similar sets 6. Interior distance on deterministic fractals}  Preprint, 1990.

\bibitem{BG}
  J. Bang-Jensen and G. Gutin, 
 \emph{Digraphs: Theory, Algorithms and Applications},
 Springer, 2009.

\bibitem{barua2022}
 A. K. Barua, S. Li, X.Li and P. Leo,
 {Asymptotic and exact self-similar evolution of growing dendrite},
 \emph{International Journal of Numerical Analysis \& Modeling}, \textbf{19}:6 (2022), 777--792.

 \bibitem{Char}
J.~J.Charatonik, W.~J.Charatonik, { Dendrites} \emph{Aportaciones Mat. Comun.}, \textbf{22} (1998), 227--253.


\bibitem{Adam}
 D. Drozdov, M. Samuel and A. Tetenov,
 {On deformation of polygonal dendrites preserving the intersection graph},
 \emph{The Art of Discrete and Applied Mathematics}, \textbf{4} (2021),  2--7.

\bibitem{DT2024}
 D. A. Drozdov and A. V. Tetenov, 
 {On the dendrite property of fractal cubes} 
 \emph{Advances in the Theory of Nonlinear Analysis and Its Application} \textbf{8}:1 (2024), 73--80.

\bibitem{TD2023}
 D. Drozdov and A. Tetenov, 
 {On the classification of fractal square dendrites}
 \emph{Advances in the Theory of Nonlinear Analysis and Its Application} \textbf{7}:3 (2023), 19--96.

\bibitem{Drozdov2025}
 D. Drozdov, 
 {Self-similar forests on fractal squares}
 \emph{Siberian Electronic Mathematical Reports} \textbf{22}:1 (2025), 385--394.

 \bibitem{Hata85}
{Hata~M.}
{\it On the structure of self-similar sets} //
{Japan Journal of Applied Mathematics, 1985.  V. 2, P. 381--414.}

 \bibitem{DH1984} A. Douady and J. H. Hubbard, \textit{Exploring the Mandelbrot set. the Orsay notes}, Publ. Math. Orsay (1984)

\bibitem{Hig}
 P. Higgins,
 \emph{Techniques of semigroup theory },
 Oxford University Press, 1992.

\bibitem{Hut}
 J. Hutchinson,
 {Fractals and self-similarity},
 \emph{Indiana Univ. Math. J.} \textbf{30}:5 (1981),  713-747.



\bibitem{Kam}
 A. Kameyama,
{Julia sets and self-similar sets},
\emph{Topology and its Applications},
\textbf{54.1-3} (1993), 241-251.

\bibitem{Kig}
 J. Kigami,
 \emph{Analysis on Fractals},
 Cambridge University Press, Number 143, 2001.

\bibitem{Kig95}
 J. Kigami,
 {Harmonic Calculus on Limits of Networks and Its Application to Dendrites},
 \emph{Journal of Functional Analysis} \textbf{128}:1 (1995), 48--86.

\bibitem{Kur}
 K. Kuratowski,
 \emph{Topology},  
 Volumes 1 and 2,  Academic Press and PWN, New York,  1966. 

\bibitem{Maz}
S. Mazurkiewicz, 
{Sur les continus homogenes}, 
\emph{Fund. Math.} \textbf{5} (1924), 137-146.


\bibitem{IFStile}
 D. Mekhontsev,
 \emph{IFStile}, software,
 2026. Available from: \url{https://ifstile.com}.


\bibitem{Men} K. Menger, Kurventheorie, Teubner Verlag, Leipzig and Berlin 1932 (reprinted by Chelsea Publ.
Co., Bronx, NY 1967).
\bibitem{Mil} E. W. Miller, Solution of the Zarankiewicz problem, Bull. Amer. Math. Soc.38 (1932), 831-834.

\bibitem{Mor}
 M. Moran,
 {Dynamical boundary of a self-similar set},
 \emph{Fundamenta Mathematicae} \textbf{160.1} (1999), 1-14.

 \bibitem{RZ2020}
 H. Rao and S.-Q. Zhang,
 {Space-filling curves of self-similar sets (III): Skeletons}
 \emph{Fractals} \textbf{28}:2 (2020), 2050028

\bibitem{RWYZ2023}
 H. Rao, Z.-Y. Wen, Q. Yuan and Y. Zhang,
 Topology automaton and conformal dimension of post-critical-finite self-similar sets,
 preprint, 2023, arXiv:2303.10320.



\bibitem{Str}
 R. S. Strichartz,
 {Isoperimetric estimates on Sierpinski gasket type fractals},
 \emph{Trans. Amer. Math. Soc.} \textbf{351} (1999), 1705–1752.


\bibitem{STV}
 M.~Samuel, A.~Tetenov and D.~Vaulin,
 Self-similar dendrites generated by polygonal systems in the plane 
 \emph{Sib.Electron.Math.Reports,}  \textbf{14}  (2017), 737–751.

\bibitem{TYK}
 A. Tetenov, I. Yudin and M. Kadirova, 
 {Finiteness properties for self-similar continua} 
 \emph{Discrete and Continuous Dynamical Systems --- Series S} (2025). doi: 10.3934/dcdss.2025057

\bibitem{W}
T. Wazewski, Sur les courbes de Jordan ne renfermant aucune courbe simple ferme.e de Jordan,
Annales de la Societe Polonaise de Mathematique 2 (1923), 49-170.

\end{thebibliography}
\noindent
Andrei Tetenov \\
Sobolev Institute of Mathematics, \\
Novosibirsk, Russia. \\
e-mail: a.tetenov@gmail.com
\bigskip
\noindent

Ivan Yudin \\
Sobolev Institute of Mathematics, \\
Novosibirsk, Russia. \\
e-mail: uivan566@gmail.com
\bigskip
\noindent

Dmitrii Drozdov \\
Sobolev Institute of Mathematics, \\
Novosibirsk, Russia. \\
e-mail: d.drozdov1@g.nsu.ru
\bigskip
\end{document}
